% ###################
% ## Filename: warehouse_intro.tex
% ## Location: D:\Project\Prej_Teaching\IE_Chula\presentation\warehouse_01_intro
% ## Created: April 27, 2011
% ## Last updated: June 7, 2011
% ## Figure Location: ..\Figures\
% ## Note: Presentation, Teaching
% ###################

%\documentclass[10pt,presentation]{beamer}
%\documentclass[10pt,handout]{beamer}
\documentclass[10pt,handout,notes=show]{beamer}


%============================================================================================================--
\usepackage{fontspec}
\usepackage{xunicode}
\usepackage{xltxtra}
\usepackage{array}
\usepackage{xcolor}
\usepackage{verbatim}
\usepackage{beamerthemesplit}
\usepackage{beamerbasetemplates}
\usepackage{graphics}
\usepackage{graphicx}
\usepackage{hyperref}
\usepackage{multirow}
\usepackage{tikz,pgfplots}
\usepackage{algorithm}
\usepackage{algorithmic}
\usepackage{soul,xcolor}
\usepackage{colortbl}
\usepackage{multicol}

\XeTeXlinebreaklocale"th"
\defaultfontfeatures{Scale=1.5, Mapping=tex-text} %{Scale=1.23}



\usetheme{Madrid}  % JuanLesPins  Rochester
           %  Berkeley  Palo Alto    with sidebar and top
           % Goettingen  Marburg     with sidebar
           %  Copenhagen  Luebeck  Warsaw

\setbeamertemplate{headline}
{%
    \begin{beamercolorbox}{section in head/foot}
        \vskip-2pt \hskip5pt
        \insertsectionnavigationhorizontal{\textwidth}{}{}
        \hskip5pt \vskip-2pt
    \end{beamercolorbox}
}%
\setbeamertemplate{footline}
{%
    \begin{beamercolorbox}{section in head/foot}
        \vskip-2pt \hskip5pt
        \insertshorttitle \hfill \insertshortdate \hskip10pt \insertframenumber / \inserttotalframenumber
        \vskip-2pt\hskip5pt
    \end{beamercolorbox}
}%
\definecolor{alertedOrange}{RGB}{255,36,0}
\setbeamercolor{alerted text}{fg=alertedOrange}
\useinnertheme{rounded}
\usecolortheme{rose}
%\setbeamercovered{transparent}
\usefonttheme{structuresmallcapsserif}
\definecolor{LRed}{rgb} {1,.8,.8}
\definecolor{LBlue}{rgb}{.8,.8,1}

%%%%%%%%%%%%%%%%%%%%%%%%%%%%%%%%%%%%%%%%%%%%%%%%%%%%%%%%%%%%%%%%%%%%%%%%%%%%%%%
% Notation:
\newtheorem{observation}{Observation}
\newcommand{\figureLocation}[1]{../figures/#1}
\newcommand{\figureLoc}[1]{../figures/#1}
\newcommand{\refLocation}[1]{../../../../Reference/#1}
\newcommand{\NextLine}{\vskip 0.05\textheight}
\newcommand{\NextHalfLine}{\vskip 0.02\textheight}
\newcommand{\RemoveLine}{\vskip -0.1\textheight}
\newcommand{\RemoveHalfLine}{\vskip -0.02\textheight}
\newcommand{\BackTab}{\hskip -0.15\textwidth}
\newcommand{\HalfBackTab}{\hskip -0.11\textwidth}
\newcommand{\DualFigureSize}{1.10\textwidth}
\newcommand{\MiniPageSize}{0.52\textwidth}
\newcommand{\code}[1]{\textcolor{blue}{#1} }
\newcommand{\rightNote}[2]{\hfill \uncover<#1>{\textbf{\code{\#\# {#2}  }}}}
\newcolumntype{P}[1]{>{\raggedright\arraybackslash}p{#1}<{}}

\author[Kittithreerapronchai]{%
  Oran~Kittithreerapronchai\inst{1} }
\institute[IE Chula]{
  \inst{1}%
  Department of Industrial Engineering, Chulalongkorn University\\
  Bangkok 10330 THAILAND \\
  }
\subject{}
\date{}

\setbeamertemplate{enumerate items}[default]
\setbeamertemplate{bibliography item}[text]
\bibliographystyle{alpha}
\newcommand{\textFont}[1]{\footnotesize{#1}}
\setcounter{MaxMatrixCols}{20}

%%%--------------------------------

\makeatletter
\defbeamertemplate{note page}{infolines}
{%
  {%
    \small
    \insertvrule{.25\paperheight}{white!90!black}
    \vskip-.25\paperheight
    \nointerlineskip
    \vbox{
      \hfill\insertslideintonotes{0.25}\hskip-\Gm@rmargin\hskip0pt%
      \vskip-0.25\paperheight%
      \nointerlineskip
      \begin{pgfpicture}{0cm}{0cm}{0cm}{0cm}
        \begin{pgflowlevelscope}{\pgftransformrotate{90}}
          {\pgftransformshift{\pgfpoint{-2cm}{0.2cm}}%
          \pgftext[base,left]{\footnotesize\the\year-\ifnum\month<10\relax0\fi\the\month-\ifnum\day<10\relax0\fi\the\day}}
        \end{pgflowlevelscope}
      \end{pgfpicture}}
    \nointerlineskip
    \vbox to .25\paperheight{\vskip0.5em
      \hbox{\insertshorttitle[width=8cm]}%
      \setbox\beamer@tempbox=\hbox{\insertsection}%
      \hbox{\ifdim\wd\beamer@tempbox>1pt{\hskip4pt\raise3pt\hbox{\vrule
            width0.4pt height7pt\vrule width 9pt
            height0.4pt}}\hskip1pt\hbox{\begin{minipage}[t]{7.5cm}\def\breakhere{}\insertsection\end{minipage}}\fi%
      }%
      \setbox\beamer@tempbox=\hbox{\insertsubsection}%
      \hbox{\ifdim\wd\beamer@tempbox>1pt{\hskip17.4pt\raise3pt\hbox{\vrule
            width0.4pt height7pt\vrule width 9pt
            height0.4pt}}\hskip1pt\hbox{\begin{minipage}[t]{7.5cm}\def\breakhere{}\insertsubsection\end{minipage}}\fi%
      }%
      \setbox\beamer@tempbox=\hbox{\insertshortframetitle}%
      \hbox{\ifdim\wd\beamer@tempbox>1pt{\hskip30.8pt\raise3pt\hbox{\vrule
            width0.4pt height7pt\vrule width 9pt
            height0.4pt}}\hskip1pt\hbox{\insertshortframetitle[width=7cm]}\fi%
      }%
      \vfil}%
  }%
  \vskip.25em
  \nointerlineskip
  %\setmainfont{TH Sarabun New}
  \begin{minipage}{\textwidth} % this is an addition
  \footnotesize{\fontspec{TH Sarabun New} \insertnote}
  \end{minipage}               % this is an addition
  %\setmainfont{Computer Modern}

}
\makeatother

\setbeamertemplate{note page}[infolines]

\title[LSCMv v2.0]{ Lecture 09: \\ Warehousing Management:\\ Concepts \& Tours \\ ~~ }
\hypersetup{
  pdftitle={09  Warehousing Management},
  pdfsubject={2104509lscm & 2184408scm handout},
  pdfproducer={XeLatex/WinEdt},
  pdfcreator={oran.k@chula.ac.th}
}

\renewcommand{\figureLocation}[1]{../../../02_warehouse/01_presentation/figures/#1}

\begin{document}
%\fontspec[BoldFont={TH Sarabun New Bold:script=thai}]{TH Sarabun New:script=thai}

%============================================================================================================--
\begin{frame}[t]
  \titlepage

  \NextLine   \NextLine
  \begin{flushright}
    \tiny{ \textbf{last updated:} \today } \hskip3ex
  \end{flushright}

  \note{
    \noindent\textbf{Purpose:} Introduce the lecture and connect it to a familiar supply-chain decision.\par
    \noindent\textbf{Leading Question:} Where have you encountered this kind of supply-chain decision?\par
    \noindent\textbf{Teaching cue:} Get one example and identify who owns the decision and what outcome they want.\par
    Key textbooks~\cite{bartholdi2009warehouse, mulcahy1994warehouse, frazelle2002world}
    \begin{itemize}
      \item welcome 2104509 Warehouse \& Warehousing management
      \item problem-based/ project-based class $\rightarrow$ know how to work
      \item theories are either very easy or very difficult
      \item lecture is designed for discussion based learning
      \item material from lecturer is English
      \item no S/U allow in this class
    \end{itemize}
  }
\end{frame}

%============================================================================================================--
\begin{frame}[t]
  \frametitle{Outline}

  \tableofcontents

  \begin{flushright}
      \tiny{ \textbf{source}: General references~\cite{bartholdi2009warehouse,mulcahy1994warehouse,frazelle2002world,oran2018warehouseENG} } \hskip3ex
  \end{flushright}

  \note{
    \noindent\textbf{Purpose:} Show how the lecture moves from concepts to applications.\par
    \noindent\textbf{Leading Question:} Which topic would you consult first for the decision raised at the opening?\par
    \noindent\textbf{Teaching cue:} Point to the relevant section and ask which question should be resolved by the end.\par
    \begin{itemize}
      \item introduction/rules
      \item important \& motivation of warehousing
      \item review an IE basic concept
    \end{itemize}
  }
\end{frame}

%============================================================================================================--
\section[Warehouse 101]{Warehousing basic concepts}
%============================================================================================================--
\begin{frame}
  \frametitle{Warehouse \& Logistics/Supply Chain}

    \begin{itemize}
      \item Warehouse $\in$ Supply Chain
      \item Warehouse $=\!\!\!?$ a bid dark, gloomy, \& messy building
      \item Warehousing management $=\!\!\!?$ Inventory management
      \item Warehouse $=\!\!\!?$ DC $=\!\!\!?$ Transit facility $=\!\!\!?$ Silo $=\!\!\!?$ Crossdock
      \item If goals of Supply Chain is to ensure that \alert{customers} got
      \begin{itemize}
      \item[] the right \alert{item}
      \item[] in the right \alert{quantity}
      \item[] at the right \alert{place}
      \item[] at the right \alert{time}
      \item[] in the right \alert{condition}
      \item[] at the right \alert{price}
      \item[] at the optimum cost to the \alert{organization(s)}
      \end{itemize}
      \item[] then, how these related to \alert{warehouse}?
    \end{itemize}

  \note{
    \noindent\textbf{Purpose:} Connect Warehouse \& Logistics/Supply Chain to a decision and the evidence needed to explain it.\par
    \noindent\textbf{Leading Question:} Which warehouse activities support the broader logistics promise?\par
    \noindent\textbf{Teaching cue:} Trace one order through inventory, handling, and dispatch.\par
    What is exactly Supply Chain? How it different from logistics?
    \begin{itemize}
      \item \textbf{Supply Chain:} typical thinking of transportation location planning \& inventory.
      \item \textbf{Item/Quantities:} otherwise customer doesn't care (e.g, chemical compound; book)
      \item \textbf{Right time/condition/price} why use 'right' $\rightarrow$ at what consumer percept (Chain 1st Quality, Intel, Harddrive)
      \item \textbf{Inventory Mgt} VS \textbf{Warehousing Mgt}
      \item \textbf{Time measurement:}
      \begin{description}
        \item[year] repeated student
        \item[day] editor
        \item[second] accident
      \end{description}
      \item \textbf{Organization:} target customers; business strategy
    \end{itemize}
  }
\end{frame}

%============================================================================================================--
\renewcommand{\textFont}[1]{\footnotesize{#1}}
\usetikzlibrary{shapes.arrows}
\tikzstyle{path}=[rectangle]
\tikzstyle{strColor}=[fill=gray!30]
\tikzstyle{ordColor}=[fill=green!30]
\tikzstyle{vasColor}=[fill=yellow!20]
\tikzstyle{mfgColor}=[fill=blue!30]
%================================================================================================-
\begin{frame}
  \frametitle{Development of Warehousing Management}

   \begin{center}
    \begin{tikzpicture}[xscale=0.5,yscale=0.3]

      \draw [strColor,rounded corners=10pt] (0,0) rectangle (3.5,10);
      \node [] (whText) at (1.5,11.5) {\begin{tabular}{c} ~ \\ \small{\textbf{Warehouse}} \end{tabular}};
      \node [] (strTextWH) at (2.0,5) {\textFont{storage}};

      \draw [strColor,rounded corners=5pt] (5,0) rectangle (8.5,5);
      \draw [ordColor,rounded corners=5pt] (5,5) rectangle (8.5,10);
      \node [] (dcText) at (6.5,11.5) {\begin{tabular}{c} \small{\textbf{Distribution}} \\ \normalsize{\textbf{Center}} \end{tabular}};
      \node [] (strTextDC) at (7.0,2.5) {\textFont{storage}};
      \node [] (ordTextDC) at (7.0,7.5) {\textFont{distribute}};

      \draw [strColor,rounded corners=3.5pt] (10,0) rectangle (13.5,3.5);
      \draw [ordColor,rounded corners=3.0pt] (10,3.5) rectangle (13.5,6.5);
      \draw [vasColor,rounded corners=3.5pt] (10,6.5) rectangle (13.5,10);
      \node [] (lcText) at (11.5,11.5) {\begin{tabular}{c} \small{\textbf{Logistics}} \\ \normalsize{\textbf{Center}} \end{tabular}};
      \node [] (strTextLC) at (12.0,2) {\textFont{storage}};
      \node [] (ordTextLC) at (12.0,5) {\textFont{distribute}};
      \node [] (vasTextLC) at (12.0,8) {\textFont{VAS}};

      \draw [strColor,rounded corners=2pt] (15,0)   rectangle (18.5,1.5);
      \draw [ordColor,,rounded corners=3pt] (15,1.5) rectangle (18.5,4.5);
      \draw [vasColor,rounded corners=3pt] (15,4.5) rectangle (18.5,8.0);
      \draw [mfgColor,rounded corners=2pt] (15,8.0) rectangle (18.5,10);
      \node [] (fcText) at (16.5,11.5) {\begin{tabular}{c} \small{\textbf{Fulfilment}} \\ \normalsize{\textbf{Center}} \end{tabular}} ;
      \node [] (strTextFC) at (17.0,0.75) {\textFont{storage}};
      \node [] (ordTextFC) at (17.0,3.00) {\textFont{distribute}};
      \node [] (vasTextFC) at (17.0,6.75) {\textFont{VAS}};
      \node [] (mfgTextFC) at (17.0,9.00) {\textFont{assembly/manuf}};

      \draw[->,black,thick] (-0.3,-0.3) -- (19.0,-0.3) node[below] {\textFont{\textbf{ Time }} };
      \draw[->,black,thick] (-0.3,-0.3) -- (-0.3,11.0) node[left,rotate=90,yshift=6.0]  {\textFont{\textbf{Portion }} };

      \node [] (myInfo) at (-.5,-2) {\begin{tabular}{r} ~ \\ \textFont{\textbf{A.D.:}} \\ \textbf{\textFont{Equipment: }} \\ \textbf{\textFont{Cycle: }} \end{tabular}};
      \node [] (whInfo) at (2.5,-2) {\begin{tabular}{c} ~ \\ \textFont{\textbf{< 1960}} \\ \textFont{ mainframe } \\ \textFont{ mth } \end{tabular}};


      \node [] (dcInfo) at (7.5,-2) {\begin{tabular}{c} ~ \\ \textFont{\textbf{1960-1980}} \\ \textFont{ PC } \\ \textFont{ week } \end{tabular}};

      \node [] (lcInfo) at (12.5,-2) {\begin{tabular}{c} ~ \\ \textFont{\textbf{1980-2000}} \\ \textFont{ internet } \\ \textFont{ day }\end{tabular}};

      \node [] (fcInfo) at (17.0,-2) {\begin{tabular}{c} ~ \\ \textFont{\textbf{2000-Present}} \\ \textFont{ mobile } \\ \textFont{ hour } \end{tabular}};

      \draw[->,black,thick] (3.5,5.0) -- (5.0,2.5) ;
      \draw[->,black,thick] (8.5,2.5) -- (10.,1.75) ;
      \draw[->,black,thick] (13.5,1.75) -- (15.,1) ;
    \end{tikzpicture}
   \end{center}



  \begin{flushright}
    \tiny{\textbf{source:} Frazelle, E. 2002. ``World-Class Warehousing \& Material Handling''~ \cite{frazelle2002world} } \hskip10em
  \end{flushright}

  \note{
    \noindent\textbf{Purpose:} Show reference\par
    \noindent\textbf{Leading Question:} What operating change distinguishes historical and newer warehouses?\par
    \noindent\textbf{Teaching cue:} Compare a documented change in volume, assortment, traceability, or service requirements.\par
%   \begin{center}
%    %\includegraphics[width=0.35\linewidth]{\figureLocation{scmTimeline.png}}
%   \end{center}

    \begin{itemize}
      \item \textbf{traffic:} manual/ many carries using \textsl{mainframe}
      \item \textbf{logistic:} max profit/low cost + measure / few carriers/ EDI/ bar code using \textsl{PC}
      \item \textbf{supply chain:} value (speed/scale/outsoucing ) / connected WMS/ using \textsl{ERP}
      \item \textbf{collaboration:} spring new business (data) \textsl{SOA/Clound}
    \end{itemize}
  }
\end{frame}

%============================================================================================================--
\begin{frame}
  \frametitle{Warehouse \textbf{VS} Warehousing management}

  \begin{description}
    \item[Warehouse] a physical location store inventory. Types of warehouse are:
      \begin{itemize}
        \item \textbf{Products:} finish goods, work-in-process, raw materials
        \item \textbf{Material handling:} unit load, break-bulk, chilled, chemical
        \item \textbf{Interactions:} picker-move (goods-to-men), picker-fixed (men-to-goods)
        \item \textbf{Business:} retail, service parts, 3PL, perishables \NextLine
      \end{itemize}
  \end{description}

   \note{
    \noindent\textbf{Purpose:} Connect Warehouse VS Warehousing management to a decision and the evidence needed to explain it.\par
    \noindent\textbf{Leading Question:} How do the physical facility and its management decisions differ?\par
    \noindent\textbf{Teaching cue:} Separate equipment and space from policies about labor, replenishment, and service.\par
    Warehouse can be anywhere \& any locations
    \begin{itemize}
      \item \textbf{WIP:} textile to reduce lead time
      \item \textbf{retail VS service part:} type of product, \# SKU \& quantity,
      \item \textbf{picker-move:} picker-2-part
      \item \textbf{product-move:} part-2-picker
    \end{itemize}

    Warehousing Mgt: oversee
    \begin{itemize}
      \item \textbf{Resources:} time, money, workforce, location, area/space
      \item \textbf{Operation:} Value-added, Non-Value added
      \item \textbf{maintaining:} no surprise
      \item \textbf{monitoring VS securing:} design; hazard
    \end{itemize}
  }
\end{frame}

%%============================================================================================================--
\begin{frame}
  \frametitle{Objectives warehousing management}

    \textsc{Objectives}
    \begin{itemize}
      \item To ensure availability of resources for \alert{planned level} of business.
      \item To meet \alert{throughput requirements}.
      \item To provide an cost effective service while meet business objectives. \NextLine
    \end{itemize}

    \textsc{Specifically:} Time, Space, \& Cost
    \begin{itemize}
      \item minimizing \alert{frequency/distance} of movement
      \item maximizing the \alert{use of cubic space}
      \item enabling the use of \alert{standard storage} \& \alert{handling equipment}
      \item \alert{speeding up} loading \& unloading
      \item minimizing damages \& thieving
    \end{itemize}


  \note{
    \noindent\textbf{Purpose:} Connect Objectives warehousing management to a decision and the evidence needed to explain it.\par
    \noindent\textbf{Leading Question:} Which goals must be balanced when improving the warehouse?\par
    \noindent\textbf{Teaching cue:} Pair throughput and accuracy with labor, space, and safety constraints.\par
    Warehouse has different goals; some are conflict at each other $\rightarrow$ what is your company trade-off
    \begin{itemize}
      \item \textbf{Third Obj:} productivity / warehouse as support unit $\rightarrow$ cost saving
      \item \textbf{Second Obj:}design/control / Warehouse $\neq$ Storage room
      \item \textbf{First Obj:} efficiency / resource must meet plan
      \item[] ลดระยะทางในการเคลื่อนย้ายให้มากที่สุด
      \item[] ใช้พื้นที่และปริมาตรในการจัดเก็บให้เกิดประโยชน์สูงสุด
      \item[] มั่นใจว่าแรงงาน เครื่องมือ อุปกรณ์ สาธารณูปโภคต่างๆ มีเพียงพอ
      \item[] สร้างความพึงพอใจในการทำงานให้แก่ผู้เกี่ยวข้อง
      \item[] ควบคุม วางแผนได้อย่างต่อเนื่อง และรักษาระดับการใช้ทรัพยากรต่างๆ
    \end{itemize}
  }
\end{frame}

%============================================================================================================-
% Standardize how tikz displays vendors, crossdock, \& stores
\usetikzlibrary{shapes,arrows,positioning,decorations.pathmorphing}
\renewcommand{\textFont}[1]{\footnotesize{#1}}
\tikzstyle{possLink}=[gray,->]
\tikzstyle{mainLink}=[black,thick,->]
\tikzstyle{block}=[rectangle,draw=black!80,thick,fill=blue!15,text centered,text width=20mm,text height=2mm]
\tikzstyle{ven}=[rectangle,draw=black!80,thick,fill=blue!15,text centered,minimum height=6mm,minimum width=6mm]
\tikzstyle{wh} =[isosceles triangle,shape border rotate=-90,isosceles triangle apex angle=60,draw=black!80,thick,fill=orange!15,text centered,minimum height=6mm]
\tikzstyle{sto}=[circle,draw=black!80,thick,fill=green!15,text centered,minimum height=6mm]
%============================================================================================================--
\begin{frame}
  \frametitle{Why do we need warehouse?}
  \begin{itemize}
    \item To prevent against \alert{fluctuations} from suppliers and/or customers (Wal-Mart, SCG dealer)
    \item To exploit \alert{economy of scale} \& \alert{fright consolidation} (THD)
    \item To perform \alert{value-added activities} (e.g. HP DeskJet, NY )
  \end{itemize}


  \begin{center}
  \begin{tikzpicture}[xscale=0.8,yscale=0.8,>=stealth]

    \node[ven](venA3)     at (0.0,4.0) {\tiny{$V_3$}};
    \node[ven](venA2)     at (0.0,2.0) {\tiny{$V_2$}};
    \node[ven](venA1)     at (0.0,0.0) {\tiny{$V_1$}};
    %\node[wh](whA)        at (1.5,2.0) {\small{$DC$}};
    \node[sto](stoA1)     at (3.0,0.0) {\tiny{$S_1$}};
    \node[sto](stoA2)     at (3.0,2.0) {\tiny{$S_2$}};
    \node[sto](stoA3)     at (3.0,4.0) {\tiny{$S_3$}};

    \draw[->,thin] (venA1.east) -- (stoA1.west);
    \draw[->,thin] (venA1.east) -- (stoA2.west);
    \draw[->,thin] (venA1.east) -- (stoA3.west);
    \draw[->,thin] (venA2.east) -- (stoA1.west);
    \draw[->,thin] (venA2.east) -- (stoA2.west);
    \draw[->,thin] (venA2.east) -- (stoA3.west);
    \draw[->,thin] (venA3.east) -- (stoA1.west);
    \draw[->,thin] (venA3.east) -- (stoA2.west);
    \draw[->,thin] (venA3.east) -- (stoA3.west);


%    \node[inner sep=0pt] at (1.50,0.25){\includegraphics[width=18pt]{\figureLocation{truck.png}}};
%    \node[inner sep=0pt] at (1.50,4.25){\includegraphics[width=18pt]{\figureLocation{truck.png}}};

    \node[ven](venB3)     at (6.0,4.0) {\tiny{$V_3$}};
    \node[ven](venB2)     at (6.0,2.0) {\tiny{$V_2$}};
    \node[ven](venB1)     at (6.0,0.0) {\tiny{$V_1$}};
    \node[wh](whB)        at (7.5,2.0) {\tiny{~}};
    \node[sto](stoB1)     at (9.0,0.0) {\tiny{$S_1$}};
    \node[sto](stoB2)     at (9.0,2.0) {\tiny{$S_2$}};
    \node[sto](stoB3)     at (9.0,4.0) {\tiny{$S_3$}};

    \draw[->,thick] (venB1.east) -- (whB.west);
    \draw[->,thick] (venB2.east) -- (whB.west);
    \draw[->,thick] (venB3.east) -- (whB.west);
    \draw[->,thick] (whB.east) -- (stoB1.west);
    \draw[->,thick] (whB.east) -- (stoB2.west);
    \draw[->,thick] (whB.east) -- (stoB3.west);

  \end{tikzpicture}
  \end{center}


    \note{
    \noindent\textbf{Purpose:} Connect Why do we need warehouse to a decision and the evidence needed to explain it.\par
    \noindent\textbf{Leading Question:} When does holding stock at a warehouse improve supply-chain service?\par
    \noindent\textbf{Teaching cue:} Compare consolidation and availability gains with added inventory and handling.\par
    40-45\% Thailand Logistics cost = inventory (Warehouse, Opportunity, Storage)
    \begin{itemize}
      \item \textbf{Cost center:} who will invest in Warehoue, no real profit, unlike MC
      \item \textbf{Ownership:} focus on core business; DHL, UPS, CEVA are thriving
      \item \textbf{Pressure:} Lux soup factory example (batch production \& retail quantity)
      \item \textbf{Entropy:} making contact or having activity $\rightarrow$ gradually mess up
      \item \textbf{Fluctuation:} Wal-mart story of bulk purchasing, SCG dealer price increases
      \item \textbf{Freight:} THD crossdock; 3PL
      \item \textbf{Value-added:} Tylenel from KOR
      \item[] ป้องกันความไม่แน่นอนของอุปสงค์ Demand และ อุปทาน Supply
      \item[] เพื่อประโยชน์ขนาดของเศรษฐกิจ Economy of scale และ Flight consolidation
      \item[] สร้างมูลค่าเพิ่มทางโลจิสติกส์ ประกอบ (HP DeskJet), ติดฉลาก (NY Pharmaceutical Law)
    \end{itemize}
  }
\end{frame}

%============================================================================================================--
\usetikzlibrary{shapes,arrows,positioning,decorations.pathmorphing}
%============================================================================================================--
\begin{frame}
  \frametitle{Units in warehouse}

   \begin{center}
    \includegraphics[width=0.40\linewidth]{\figureLocation{01intro/unitInWarehouse.jpg}}
   \end{center}

  \begin{flushright}
    \tiny{\textbf{source:} Bartholdi, J. \& Hackmans, S. 2009.~\cite{bartholdi2009warehouse}\hskip10pt}
  \end{flushright}

  \begin{alertblock}{}{}{}
    \textbf{Stock Keeping Unit (SKU)} a unit identification (consider difference) in a warehouse. a \alert{pallet} of red pens $\neq$ a \alert{single} red pen
  \end{alertblock}

  \note{
    \noindent\textbf{Purpose:} Connect Units in warehouse to a decision and the evidence needed to explain it.\par
    \noindent\textbf{Leading Question:} Why must the warehouse keep a consistent unit of measure?\par
    \noindent\textbf{Teaching cue:} Convert an example from each to case to pallet and check the resulting inventory quantity.\par
    What is the biggest unitload possible in Supply Chain?
    \begin{itemize}
      \item \textbf{Pallet:} with shrinkage wrap $\rightarrow$ 20 years before adaptation in Thailand
      \item \textbf{Tier:} sub-pallet (กระเบื้องลอนคู่ -- 10 แผ่น)
      \item \textbf{Case:} found in stores
      \item \textbf{Inner pack:} consumer
      \item \textbf{Unit:} consumer
    \end{itemize}
  Not all products have all units (e.g., กระเบื้อง มีเฉพาะ pallet-case-unit)
  }
  \end{frame}

%============================================================================================================--
\begin{frame}
  \frametitle{Warehouse as flow process system}

  \begin{center}
  \begin{tikzpicture}[xscale=1,yscale=1,>=stealth]
    %\draw[step=.2,gray,very thin]  (-2.0,-2.0)  grid  (2.0,2.0);
    \draw[fill=gray!25] (-3,-0.4)   ellipse (0.1 and 0.4) ;
    \draw (-3,0.8)   arc (90:270:0.1 and 0.4) ;
    \draw (-3,-0.8) -- ++(0:1.5) -- ++ (90:1.6) -- ++ (180:1.5);
    \draw[->,blue,thick] (-3.5,-0.25) -- ++(0:0.5) ;
    \draw[->,blue,thick] (-3.5,0.00) -- ++(0:0.5) ;
    \draw[->,blue,thick] (-3.5,0.25) -- ++(0:0.5) ;

    \draw (-1.5,0.0) -- ++(90:1.5) -- ++ (0:3) -- ++ (-90:3) -- ++(180:3) -- ++ (90:1.5);
    \draw[-,blue,thick,decorate,decoration=snake] (-1.5,-0.5) --  (1.5,-0.5)  ;

    \draw[fill=gray!25] (3,-0.2)    ellipse (0.05 and 0.2) ;
    \draw (3,0.0)  arc (-90:90:0.05 and 0.2) ;
    \draw (3,-0.4) -- ++(180:1.5) -- ++ (90:0.8) -- ++ (0:1.5);
    \draw[->,blue,thick] (3.00,0.0) -- ++(0:1.5) ;

    \node(incoming)   at (-2.5,-2.0) {\small{Incoming flow}};
    \node(incoming)   at (2.5,-2.0) {\small{Outgoing flow}};
    \node(incoming)   at (0.0,-1.0) {\small{Warehouse}};
  \end{tikzpicture}
  \end{center}

  \begin{itemize}
    \item \textbf{Water:} single SKU, and compressible flow
    \item \textbf{Flow balancing:} \alert{incoming} flow \& \alert{outgoing} flow
    \item \textbf{Keep flow moving:} avoid double handling; space blocking
    \item \textbf{Smooth Flow:} resolve bottlenecks, avoid layouts that impede smooth
  \end{itemize}

  \note{
    \noindent\textbf{Purpose:} Intro flow in warehouse\par
    \noindent\textbf{Leading Question:} What limits throughput in this warehouse flow?\par
    \noindent\textbf{Teaching cue:} Mark queues at receiving, storage, picking, and shipping and choose one stage to measure.\par
    \begin{itemize}
      \item pressure $\rightarrow$ picking frequency
      \item diameter $\rightarrow$ volume of products
      \item \textbf{In textbook:}Warehouse as System (queuing system)
      \item If system is in Steady Stage, $\exists$ balance between inflow \& outflow
    \end{itemize}
  }
  \end{frame}

%============================================================================================================--
\begin{frame}
  \frametitle{Components in warehouse}

  \begin{itemize}
      \item \textbf{Facility:} building, yard, surrounding
      \item \textbf{Human:} manager, picker, checker, IT, consult
      \item \textbf{Material Handling:} products, \alert{storage location}, equipments
      \item \textbf{Processes:} \alert{main activities}, value-added logistics (VAL), counting, reconcile, document
  \end{itemize}

  \NextHalfLine
  \textsc{Principle in warehousing management}
    \begin{itemize}
      \item \textbf{F.A.S.T:}
      \begin{itemize}
        \item \textbf{Flow}: minimizing total movements/cost \alert<2->{ $\rightarrow$ no double handling}
        \item \textbf{Accessibility}: inside \& outside buildings \alert<2->{ $\rightarrow$ no blockage}
        \item \textbf{Space}: $\approx$ 40\% of cost related to warehouse; \alert<2->{ $\rightarrow$ use height}
        \item \textbf{Throughput}: equipments \alert<2->{ $\rightarrow$ use right \& flexible ones}
      \end{itemize}
      \item \textbf{Planning:} long-term goal \& short-term req$^m$, control \& feedback,
      %\item \textbf{Standardization:} standardized equipments, material handling, flow, workload
      \item \textbf{House keeping:} clean, neat, safety, security \& eco-friendly
      \item \textbf{Flexibility:} free space, stacking area, multi-purpose equipment
    \end{itemize}

    \note{
    \noindent\textbf{Purpose:} Connect Components in warehouse to a decision and the evidence needed to explain it.\par
    \noindent\textbf{Leading Question:} How do flow, accessibility, space, and throughput trade off?\par
    \noindent\textbf{Teaching cue:} Compare higher storage density with accessibility and safe movement for a fast-moving item.\par
    \begin{itemize}
      \item \textbf{Unitization:} IE mindset; postponement (group together, differentiate as late as possible)
      \item \textbf{House keeping:} typical IE work
      \item \textbf{Flexible:} buffer area; small very connivent; tough rules
      \item \textbf{Flex VS Unitization:} Std work, but Flex resource
      \item \textbf{Ergonomic:} eliminate repletion (mental \& physical), hard labor, work place
    \end{itemize}

    \begin{itemize}
      \item \textbf{Facility:} WH for JIT is near plant, WH for consumer product near BKK
      \item \textbf{Human:} discipline, nationality, motivation-incentive
      \item \textbf{Material Handling$^{*}$:} products (nature \& law GMP), storage \& equipments (calling \& trade-off)
      \item \textbf{Processes$^{*}$:} focuses
    \end{itemize}
  }
  \end{frame}

%============================================================================================================--
\begin{frame}
  \frametitle{Units in warehouse}

   \begin{center}
    \includegraphics[width=0.40\linewidth]{\figureLocation{01intro/unitInWarehouse.jpg}}
   \end{center}

  \begin{flushright}
    \tiny{\textbf{source:} Bartholdi, J. \& Hackmans, S. 2009.~\cite{bartholdi2009warehouse}\hskip10pt}
  \end{flushright}

  \begin{alertblock}{}{}{}
    \textbf{Stock Keeping Unit (SKU)} a unit identification (consider difference) in a warehouse. a \alert{pallet} of red pens $\neq$ a \alert{single} red pen
  \end{alertblock}

  \note{
    \noindent\textbf{Purpose:} Connect Units in warehouse to a decision and the evidence needed to explain it.\par
    \noindent\textbf{Leading Question:} Which handling unit should be chosen for this order profile?\par
    \noindent\textbf{Teaching cue:} Compare pieces, cases, and pallets for handling effort and service.\par
    What is the biggest unitload possible in Supply Chain?
    \begin{itemize}
      \item \textbf{Pallet:} with shrinkage wrap $\rightarrow$ 20 years before adaptation in Thailand
      \item \textbf{Tier:} sub-pallet (กระเบื้องลอนคู่ -- 10 แผ่น)
      \item \textbf{Case:} found in stores
      \item \textbf{Inner pack:} consumer
      \item \textbf{Unit:} consumer
    \end{itemize}
  Not all products have all units (e.g., กระเบื้อง มีเฉพาะ pallet-case-unit)
  }
  \end{frame}

%============================================================================================================--
\usetikzlibrary{shapes,arrows,positioning,decorations.pathmorphing}
%============================================================================================================--
\begin{frame}
  \frametitle{Warehouse as flow process system}

  \begin{center}
  \begin{tikzpicture}[xscale=1,yscale=1,>=stealth]
    %\draw[step=.2,gray,very thin]  (-2.0,-2.0)  grid  (2.0,2.0);
    \draw[fill=gray!25] (-3,-0.4)   ellipse (0.1 and 0.4) ;
    \draw (-3,0.8)   arc (90:270:0.1 and 0.4) ;
    \draw (-3,-0.8) -- ++(0:1.5) -- ++ (90:1.6) -- ++ (180:1.5);
    \draw[->,blue,thick] (-3.5,-0.25) -- ++(0:0.5) ;
    \draw[->,blue,thick] (-3.5,0.00) -- ++(0:0.5) ;
    \draw[->,blue,thick] (-3.5,0.25) -- ++(0:0.5) ;

    \draw (-1.5,0.0) -- ++(90:1.5) -- ++ (0:3) -- ++ (-90:3) -- ++(180:3) -- ++ (90:1.5) ;
    \draw[-,blue,thick,decorate,decoration=snake] (-1.5,-0.5) --  (1.5,-0.5)  ;

    \draw[fill=gray!25] (3,-0.2)    ellipse (0.05 and 0.2) ;
   \draw (3,0.0)  arc (-90:90:0.05 and 0.2) ;
    \draw (3,-0.4) -- ++(180:1.5) -- ++ (90:0.8) -- ++ (0:1.5);
    \draw[->,blue,thick] (3.00,0.0) -- ++(0:1.5) ;

    \node(incoming)   at (-2.5,-2.0) {\small{Incoming flow}};
    \node(incoming)   at (2.5,-2.0) {\small{Outgoing flow}};
    \node(incoming)   at (0.0,-1.0) {\small{Warehouse}};
  \end{tikzpicture}
  \end{center}

  \begin{itemize}
    \item \textbf{Flow balancing:} \alert{incoming} flow \& \alert{outgoing} flow
    \item \textbf{Keep flow moving:} avoid double handling; space blocking
    \item \textbf{Smooth Flow:} resolve bottlenecks, avoid layouts that impede smooth
  \end{itemize}

  \note{
    \noindent\textbf{Purpose:} Connect Warehouse as flow process system to a decision and the evidence needed to explain it.\par
    \noindent\textbf{Leading Question:} Where should dwell time be measured in the process?\par
    \noindent\textbf{Teaching cue:} Mark receipt and dispatch timestamps and identify what causes waiting between them.\par
    \begin{itemize}
      \item pressure $\rightarrow$ picking frequency
      \item diameter $\rightarrow$ volume of products
      \item \textbf{In textbook:}Warehouse as System (queuing system)
      \item If system is in Steady Stage, $\exists$ balance between inflow \& outflow
      \item[] Balancing flows: สินค้าเข้า = สินค้าออก
      \item[] Smooth flows: กำจัด double handling และ bottleneck
    \end{itemize}
  }
  \end{frame}

%============================================================================================================--
\begin{frame}
  \frametitle{Common Classification of Warehouse}

  \begin{description}
    \item[Warehouse] a physical location store inventory. Types of warehouse, providing useful \alert{insight}, are:
      \small{
      \begin{itemize}
        \item \textbf{Products:} finish goods, work-in-process, raw materials
        \item \textbf{Unit of handling:}  MHE: pallet, carton, piece
        \item \textbf{Nature of Storage:} security, storage requirement
        \item \textbf{Storage Policy:} dedicated $\leftarrow$ class-based $\rightarrow$ shared/random
        \item \textbf{Temperature:} Frozen $\leftarrow$ Chilled $-$ Air Condition $\rightarrow$ Ambiance
        \item \textbf{Management:} Public warehouse $\leftrightarrow$ Private warehouse
        \item \textbf{Movement:} Men-to-Goods $\leftarrow$ Goods-to-Men $\rightarrow$ Automation
        \item \textbf{Business:} retail, service parts, 3PL (DC), fulfillment
      \end{itemize} }
  \end{description}


  \note{
    \noindent\textbf{Purpose:} Connect Common Classification of Warehouse to a decision and the evidence needed to explain it.\par
    \noindent\textbf{Leading Question:} How does the function of a warehouse shape its layout?\par
    \noindent\textbf{Teaching cue:} Compare storage, distribution, and cross-docking on dwell time and activity mix.\par
    \begin{itemize}
      \item $<<$empty$>>$
    \end{itemize}
  }

\end{frame}

%============================================================================================================--
\begin{frame}
  \frametitle{Example: 7Eleven store as a warehouse}

  \begin{center}
    \includegraphics [height=0.45\textheight]{\figureLocation{04tours/7eleven.png}}
  \end{center}

  \small{
  \NextHalfLine
  \begin{tabular}{lr}
    \multicolumn{2}{l}{$\bullet$ \textbf{Unit of handling:}  FG piece picking (basket) }\\
    \multicolumn{2}{l}{$\bullet$ \textbf{Nature of Storage:} security and authorize (back counter)}\\
    \multicolumn{2}{l}{$\bullet$ \textbf{Temperature:} frozen (ice cream,  AirCon (Milk, Drink), ambiance (Grocery)}    \\
    $\bullet$  \textbf{Policy:} dedicate-class & $\bullet$ \textbf{Management:} private warehouse\\
    $\bullet$ \textbf{Movement:} Men-to-Goods  & \\
    \multicolumn{2}{l}{$\bullet$ \textbf{Others (TBR):} 24/7, U-shaped layout, vertical bin-shelf racking} \\
  \end{tabular}
  }

  \note{
    \noindent\textbf{Purpose:} Connect Example: 7Eleven store as a warehouse to a decision and the evidence needed to explain it.\par
    \noindent\textbf{Leading Question:} Which warehouse functions are visible in a convenience store?\par
    \noindent\textbf{Teaching cue:} Trace replenishment, storage, picking, and inventory counting for one product.\par
    \begin{itemize}
      \item $<<$empty$>>$
    \end{itemize}
  }

  \end{frame}

%============================================================================================================--
\begin{frame}
  \frametitle{Inside CP All DC (7Eleven warehouse) }

  \begin{center}
    \includegraphics [height=0.55\textheight]{\figureLocation{04tours/cpall.png}}
  \end{center}

  \begin{itemize}
    \item \textbf{Nature:} distributing products in 7Eleven for \alert{DC Fee} and QC
    \item \textbf{Receiving:} 10w or 18w supplier delivery as carton or pallet
    \item \textbf{Equipment:} trolley, tote, RT, pick-to-light
    \item \textbf{Picking:} wave picking (heavy, pieces $\rightarrow$ tote), chilled
    \item \textbf{Shipping:} 4w outsource as tote + beverage
  \end{itemize}


  \note{
    \noindent\textbf{Purpose:} Connect Inside CP All DC (7Eleven warehouse) to a decision and the evidence needed to explain it.\par
    \noindent\textbf{Leading Question:} What throughput or service decision is visible in the distribution center case?\par
    \noindent\textbf{Teaching cue:} Separate observations from site-specific claims and identify additional evidence to collect.\par
    \begin{itemize}
      \item $<<$empty$>>$
    \end{itemize}
  }

  \end{frame}

%============================================================================================================--
\section[Activity]{Basic concept of warehousing activities }
%================================================================================================-
\begin{frame}
  \frametitle{Warehousing activities}

%------------------------------------------------------
\tikzstyle{block}=[rectangle,draw=black!80,thick,fill=darkgray!15,text centered,text width=16mm,text height=1.0mm]
\tikzstyle{blockEx}=[rectangle,draw=black!80,thick,fill=orange!15,text centered,text width=16mm,text height=1.0mm]
\tikzstyle{blockIn}=[rectangle,draw=black!80,thick,fill=blue!15,text centered,text width=16mm,text height=1.0mm]
\tikzstyle{blockOut}=[rectangle,draw=black!80,thick,fill=magenta!15,text centered,text width=16mm,text height=1.0mm]
 \begin{center}
   \begin{tikzpicture}[xscale=0.7,yscale=0.7,>=stealth]

    \node[blockEx](supp)       at (0.0,0.0) {\textFont{Support \\ Function}};
    \node[blockIn](put)        at (0.0,1.5) {\textFont{Put-away}};
    \node[block](stor)         at (3.0,0.8) {\textFont{~\\ ~\\ \textbf{Storage} \\ ~\\ ~\\ }};
    \node[blockOut](pick)      at (3.0,-1.5) {\textFont{Picking}};
    \node[blockOut](sort)      at (0.0,-1.5) {\textFont{Sorting}};
    \node[blockOut](ship)      at (-3.0,-1.5) {\textFont{Shipping}};
    \node[blockEx](cros)       at (-3.0,0.0) {\textFont{Cross-docking}};
    \node[blockIn](reci)       at (-3.0,1.5) {\textFont{Receiving}};

    \node[inner sep=0pt] (inTruck)  at (-6.50,1.5){\includegraphics[width=40pt]{\figureLocation{trucks.png}}};
    \node[inner sep=0pt] (outTruck) at (-6.50,-1.5){\includegraphics[width=40pt]{\figureLocation{trucks.png}}};

    \draw[brown,-,thick] (-4.5,2.25) -- (4.5,2.25) -- (4.5,-2.25) -- (-4.5,-2.25) -- (-4.5,2.25);
    \draw[black,->,thick] (inTruck.east) -- (reci.west);
    \draw[black,->,thick] (ship.west) -- (outTruck.east);

    \draw[black,->,thick] (reci) -- (put) ;
    \draw[black,->,thick] (reci) -- (cros) ;
    \draw[black,->,thick] (put)  -- (2.0,1.5) ;
    \draw[black,->,thick] (stor) -- (pick);
    \draw[black,->,thick] (pick) -- (sort);
    \draw[black,->,thick] (cros) -- (ship)  ;
    \draw[black,->,thick] (sort) -- (ship)  ;

    \draw[gray, <->,dotted,thick] (supp.north) -- (put.south) ;
    \draw[gray, <->,dotted,thick] (supp.north west) -- (reci.south east) ;
    \draw[gray, <->,dotted,thick] (supp.south west) -- (ship.north east) ;
    \draw[gray, <->,dotted,thick] (supp.south) -- (sort.north) ;

   \end{tikzpicture}
%  \includegraphics[width=0.80\textwidth]{\figureLocation{02activity/activityInWarehouse.jpg}}
 \end{center}

%------------------------------------------------------

  \begin{flushright}
    \tiny{\textbf{source:} Frazelle, E. 2001.~\cite{frazelle2002world} \hskip10pt }
  \end{flushright}


  \textsc{Warehouse activity breakdown}
  \begin{center} \small{
  \begin{tabular}{lr} \hline
    Activities        & Example  \\ \hline
    Receiving  (10\%) & yard mgt, inspection, unitization  \\
    Put-away   (15\%) & positioning, slotting, stock keeping  \\
    Picking (55\%)    & \alert{dispatching} , routing \\
    Shipping (20\%)   & sorting, \alert{loading} , checker \\ \hline
  \end{tabular}
  }
  \end{center}
  \begin{flushright}
    \tiny{\textbf{source:} Bartholdi, J. \& Hackmans, S. 2009.~\cite{bartholdi2009warehouse}\hskip10pt}
  \end{flushright}


  \note{
    \noindent\textbf{Purpose:} Connect Warehousing activities to a decision and the evidence needed to explain it.\par
    \noindent\textbf{Leading Question:} Which information should be captured at each physical handoff?\par
    \noindent\textbf{Teaching cue:} Track item, quantity, location, and timestamp through the activity flow.\par
    Physical flow is easy; how about information flow
    \begin{itemize}
      \item \textbf{Receiving:} invoice (check with P/O), bill-of-lading, weight/photo/scan,
      \item \textbf{Put-Away:} stock keeping (turn whitley)
      \item \textbf{Picking:} equipments, pick-list,
      \item \textbf{Shipping:} checker
      \item \textbf{Others:} physical counting
      \item \textbf{Crossdocking:} unloading, sorting, loading
    \end{itemize}
  }
  \end{frame}

%================================================================================================-
\begin{frame}
  \frametitle{Receiving \& Put-away activities}

  \textsc{Receiving}
  \begin{itemize}
    \item \textbf{Idea:} unloading \& preparation
    \item \textbf{Importance:} initiating all operations \& activities
    \item \textbf{Basic:} doing paper work \& checking for quantity \& quality
  \end{itemize}

  \NextLine
  \textsc{Put-away}
  \begin{itemize}
    \item \textbf{Idea:} moving SKU to 'designated' locations
    \item \textbf{Importance:} defining all works downstream.
    \item \textbf{Basic:} recording where SKUs are consistently
  \end{itemize}

  \note{
    \noindent\textbf{Purpose:} Connect Receiving \& Put-away activities to a decision and the evidence needed to explain it.\par
    \noindent\textbf{Leading Question:} Which checks are required before a receipt becomes available stock?\par
    \noindent\textbf{Teaching cue:} Compare expected and received quantity and identify the quarantine or exception path.\par
    \begin{itemize}
      \item \textbf{notification:}  for high \# shipments \& low receive capacity (Tesco)
      \item \textbf{good acceptance:} compare bill of lading with PO
      \item \textbf{good recept:} accept the good
      \begin{itemize}
        \item \code{Direct shipment:} direct shipping to consumer
        \item \code{Pre-receiving:} employ information before products arrive
        \item \code{Scheduling:} balancing workforce (hiring temp), equipments
        \item \code{Crossdock:} JIT warehouse
      \end{itemize}
      \item \textbf{mixed pallet:} small \& few requested item in single pallet
    \end{itemize}

    \begin{itemize}
      \item \textbf{storage location:} easy to find, but incorrect balance
      \begin{itemize}
        \item \code{Direct put-away:} combine receiving \& put-away (no inspection w/ RFID)
        \item \code{WMS Directed put-away:} use IT to suggest locations
        \item \code{Batched:} source products for efficient put-away
        \item \code{Inter-leaving:} combine put-away with picking  \NextHalfLine
      \end{itemize}
    \end{itemize}
  }
\end{frame}

%================================================================================================-
\begin{frame}
  \frametitle{Storage and Dwell Time}

    \begin{itemize}
      \item \textbf{Idea:} preventing SKUs for damage and/or degrading
      \item \textbf{Importance:} what other thinking of warehouse
      \item \textbf{Basic:} utilizing space, while maintain easy access
      \item \textbf{storage area:}
      \begin{itemize}
        \item \code{Forward:} storing \alert{products} for carton/case/piece picking and customer delivery
        \item \code{Reserve:} storing  \alert{pallets} for pallet picking or refill other area
      \end{itemize}
      \item \textbf{storage policy:}
      \begin{itemize}
        \item \code{Dedicated:} reserved \alert{specific space} for each SKUs
        \item \code{Shared:} \alert{no reservation} \NextHalfLine
      \end{itemize}
    \end{itemize}


  \note{
    \noindent\textbf{Purpose:} Connect Non-productive activity: Storage to a decision and the evidence needed to explain it.\par
    \noindent\textbf{Leading Question:} When is storage a useful buffer and when is it avoidable dwell?\par
    \noindent\textbf{Teaching cue:} Compare protection of service with space usage, damage, and carrying cost.\par
    If products moves swiftly, you can ignore this slide
    \begin{itemize}
      \item a
    \end{itemize}
  }
\end{frame}

%================================================================================================-
\begin{frame}
  \frametitle{Between receiving \& shipping}

  \begin{center}
    \includegraphics[width=0.65\linewidth]{\figureLocation{02activity/storageArea.jpg}}
  \end{center}

  \begin{flushright}
    \tiny{\textbf{source:} Mulcahy, D. 1994.~\cite{mulcahy1994warehouse} \hskip10pt }\\
  \end{flushright}


  \note{
    \noindent\textbf{Purpose:} Connect Between receiving \& shipping to a decision and the evidence needed to explain it.\par
    \noindent\textbf{Leading Question:} Where can a delay between receiving and shipping constrain the operation?\par
    \noindent\textbf{Teaching cue:} Identify the queue, its owner, and a measure of accumulated work.\par
    \begin{itemize}
      \item $<<$empty$>>$
    \end{itemize}

  }
\end{frame}

%================================================================================================-
\begin{frame}
  \frametitle{Picking activity}

    \begin{itemize}
      \item \textbf{Idea:} getting SKU from 'designated' locations
      \item \textbf{Importance:} majority of \alert{costs \& times} incurred
      \item \textbf{Basic:} distributing 'order' \& checking 'right' quantity
      \item \textbf{Principle:}
      \begin{itemize}
        \item<1-|handout:1> \code{Minimize loose items} \& \code{staging for shipping}
        \item<1-|handout:1> \code{Minimize paperwork \& time}
      \end{itemize}
      \item \textbf{Variation:}
      \begin{itemize}
        \item \code{Single [order]:}  one tour for one order
        \item \code{Batch [order]:} one tour for many orders
        \item \code{Zone:} set area for each picker
        \item \code{Wave:} coordinated between zones \NextHalfLine
      \end{itemize}
    \end{itemize}


  \note{
    \noindent\textbf{Purpose:} Connect Picking activity to a decision and the evidence needed to explain it.\par
    \noindent\textbf{Leading Question:} What makes picking expensive or error-prone in this design?\par
    \noindent\textbf{Teaching cue:} Break down travel, search, confirmation, and return to staging time.\par
    This is the most important topics in WH so we will have a whole lecture dedicated for it

  }
\end{frame}

%================================================================================================-
\begin{frame}
  \frametitle{Picking Variation}

  \begin{center}
    \only<1|handout:0>{
      \includegraphics[width=0.55\linewidth]{\figureLocation{02activity/SingleOrderPicking_OKC.png}}
    }
    \only<2|handout:0>{
      \includegraphics[width=0.55\linewidth]{\figureLocation{02activity/BatchPicking_OKC.png}}
    }
    \only<3|handout:0>{
      \includegraphics[width=0.55\linewidth]{\figureLocation{02activity/ZoneOrderPicking_OKC.png}}
    }
    \only<4|handout:1>{
      \includegraphics[width=0.55\linewidth]{\figureLocation{02activity/ZoneBatchPicking_OKC.png}}
    }
  \end{center}


  \note{
    \noindent\textbf{Purpose:} Connect Picking Variation to a decision and the evidence needed to explain it.\par
    \noindent\textbf{Leading Question:} Which order profiles favor batch, zone, or individual picking?\par
    \noindent\textbf{Teaching cue:} Compare order overlap, deadline, and sorting effort for a concrete queue.\par
      \begin{itemize}
        \item \code{Single order:}  no sorting; more travel
        \item \code{Batch [order]:} less travel; need to sort
        \item \code{Zone:} less travel + learning; balancing work + need combine
    \end{itemize}
  }
\end{frame}

%================================================================================================-
\begin{frame}
  \frametitle{Shipping \& Other activities}

  \textsc{Sorting \&  Shipping}
  \begin{itemize}
    \item \textbf{Idea:} preparing \& checking SKU before leaving warehouse
    \item \textbf{Importance:} define productivity of warehouse, \alert{quality control}
    \item \textbf{Basic:} checking SKU, documenting transactions \& loading in reverse order
  \end{itemize}

  \NextHalfLine
  \textsc{Others}
  \begin{itemize}
    \item \textbf{Value-added logistics:} re-boxing \& re-instruction (DKSH), sample assembly (Hefale), measuring \& cutting (UF), MUJI price labeling (CRC)
    \item \textbf{Reverse logistics:} dispose product (Tesco), cleaning tote (7-Eleven)
    \item \textbf{Non value-added logistics:}  visual inspection, counting, re-location (TUF), safety training(HomePro)
  \end{itemize}

  \note{
    \noindent\textbf{Purpose:} Connect Shipping \& Other activities to a decision and the evidence needed to explain it.\par
    \noindent\textbf{Leading Question:} What checks should confirm an order before shipment?\par
    \noindent\textbf{Teaching cue:} Require item and quantity reconciliation, dock assignment, and shipment confirmation.\par
    \begin{itemize}
      \item \textbf{Siam Kubota:} put dead stock in stack-able box
      \item \textbf{TUF:} get all batch for label $\rightarrow$ return what left
      \item \textbf{HomePro: } important safety, forklift school, retail practice/ training \& re-training
    \end{itemize}
  }
\end{frame}

%============================================================================================================--
\section[Equipments]{Equipment typically found in a Warehouse}
%============================================================================================================--
\begin{frame}
  \frametitle{Benefits of equipment}

  \begin{itemize}
    \item Reduce cost (labor + space)
    \item Enhance responsiveness
    \item Maintain \alert{qualities} of products \& operations
  \end{itemize}

  \NextLine
  \textsc{Classification of equipments}
  \begin{itemize}
    \item \textbf{Unitizing equipment:} unit load $\rightarrow$ container, pallet, tote,
    \item \textbf{Storage \& retrieval equipment:} floor stack, rack, carousel, bin-shelve
    \item \textbf{Material handling equipment:} forklift, hand truck, VNA, conveyor
    \item \textbf{Automatic identification \& communication equipment:} portable bar code reader, RFID
  \end{itemize}

  \note{
    \noindent\textbf{Purpose:} Connect Benefits of equipment to a decision and the evidence needed to explain it.\par
    \noindent\textbf{Leading Question:} What evidence would justify buying handling equipment?\par
    \noindent\textbf{Teaching cue:} Compare throughput, damage, maintenance, and training with the current process.\par
    What do you need to know about before selecting equipment?
    \small{
    \begin{itemize}
      \item \textbf{Planning:} comm btw designer \& user; existing methods; flexible to handle change;  short-term req \& long-term goal
      \item \textbf{Standardization:} equipment \& method for various type
      \item \textbf{Unit load:} less effort
      \item \textbf{System:} minimize inventory, ID system
      \item \textbf{Work:} combine work, double handling, simplified work
      \item \textbf{'Other':} ergonomic, automation, space utilization, life cycle
    \end{itemize}
    }
  }
  \end{frame}

%============================================================================================================--
\begin{frame}
  \frametitle{Concept of unitizing equipment}

  \begin{center}
    \includegraphics[width=0.45\linewidth]{\figureLocation{03equipment/unitload/wrapping.jpg}}
   \end{center}

  \begin{itemize}
    \item \textbf{Idea:} standardizing items/SKUs \& making them easy to move \& collect
    \item \textbf{Where:} supplier site, receiving \& shipping area
    \item \textbf{Issues:} installation cost,  volume,  size \& shape  (7Eleven tote, Lotus cool box)
    \item \textbf{Example:} \alert{pallet}, wrapping machine
  \end{itemize}

  \note{
    \noindent\textbf{Purpose:} Connect Concept of unitizing equipment to a decision and the evidence needed to explain it.\par
    \noindent\textbf{Leading Question:} What size or stability constraint determines the handling unit?\par
    \noindent\textbf{Teaching cue:} Match product and load dimensions to the storage and movement interfaces.\par
    Pro
    \begin{itemize}
      \item More items can be handled at the same time
      \item Standardization
    \end{itemize}
    Con
    \begin{itemize}
      \item Time spent forming \& breaking down the unit load.
      \item Cost of containers/pallets \& other load restraining materials used in the unit load
      \item Empty containers/pallets may need to be returned to their point of origin.
    \end{itemize}
    Not so trivial: unitload tends to be large \\
    collect \& move many individual items
    common as RM \& as FG
  }
  \end{frame}

%============================================================================================================--
\begin{frame}
  \frametitle{Pallet \& co }

  \begin{center}
  \begin{tabular}{cc}
    \begin{minipage}{0.40\linewidth}
      \includegraphics[width=0.85\textwidth]{\figureLocation{03equipment/unitload/plasticPalletNew.jpg}}
    \end{minipage}
    &
    \begin{minipage}{0.40\linewidth}
      \includegraphics[width=0.85\textwidth]{\figureLocation{03equipment/unitload/ironWireWoodenPalletNew.jpg}}
    \end{minipage}
    \\
    \begin{minipage}{0.40\linewidth}
      \includegraphics[width=0.85\textwidth]{\figureLocation{03equipment/unitload/tierRackPalletNew.jpg}}
    \end{minipage}
    &
    \begin{minipage}{0.40\linewidth}
      \includegraphics[width=0.85\textwidth]{\figureLocation{03equipment/unitload/truckPalletNew.jpg}}
    \end{minipage}
  \end{tabular}
  \end{center}


    \begin{itemize}
      \item<1-|handout:1> \textbf{Idea:} creating unit load by std platform
      \item<1-|handout:1> \textbf{Issues:} circulation, size, form
    \end{itemize}

  \note{
    \noindent\textbf{Purpose:} Connect Pallet \& co to a decision and the evidence needed to explain it.\par
    \noindent\textbf{Leading Question:} Which dimensions must be checked across pallets, racks, and transport equipment?\par
    \noindent\textbf{Teaching cue:} Trace the same load through forks, rack supports, and vehicle space.\par
    \begin{itemize}
      \item \textbf{pallet:} different material (wooden), size (1.2 x 1.2m), form (reversible)
      \item \textbf{circulation:} save on pallet by shipping with no pallet (handling cost \textbf{VS} cost pallet)
    \end{itemize}
    Tier rack frame/ Stacking Frame
    \begin{itemize}
      \item Interlocking units that enable stacking of a load so that crushing does not occur
      \item Can be disassembled \& stored compactly when not in used
      \item Pallet frames can be used to enable multilevel block stacking
    \end{itemize}
  }
  \end{frame}

%============================================================================================================--
\begin{frame}
  \frametitle{Comparison of pallet materials}


  \begin{center}
    \includegraphics[width=0.55\textwidth]{\figureLocation{03equipment/unitload/palletMat.png}}
  \end{center}
  \begin{flushright}
    \tiny{\textbf{source:} https://www.palltechpallets.co.uk \hskip3pt}
  \end{flushright}

  \NextLine
  \begin{center} \small{
  \begin{tabular}{lrrll} \hline
    Material	     & Durability	   & Repairable	& Env. Impact 	         & Application \\ \hline
    Wood 	         & med	         & yes 	      & recyclable 	           & common \\
    Pressed Wood   & med           & yes	      & recyclable             & printing, tiber  \\
    Fiberboard	   & low 	         & no 	      & recyclable             & paper, garment \\
    Plastic 	     & high 	       & no	        & closed loop 	         & cement, automotive \\
    Metal	         & high 	       & depends 	  & closed loop 	         & grocery, food, military \\ \hline
  \end{tabular}
  }
  \end{center}



  \note{
    \noindent\textbf{Purpose:} type of pallet\par
    \noindent\textbf{Leading Question:} What conditions determine whether a pallet material is a good choice?\par
    \noindent\textbf{Teaching cue:} Compare loads, sanitation, repair, reuse cycles, and supplier specifications.\par
    \begin{itemize}
      \item \textbf{พาเลทไม้} เป็นวัสดุที่นิยม มีราคาไม่แพงนัก สามารถนำกลับมาใหม่ได้โดยการซ่อม
      \item \textbf{พาเลทโลหะ} เป็นพาเลทภายในมีความทนทานสูงแต่มีราคาสูงและน้ำหนักมาก อายุการใช้งานประมาณ 5-8 ปี ใช้เก็บสินค้าที่มีน้ำหนักมาก
      \item \textbf{พาเลทพลาสติก} เป็นพาเลทภายในมีความทนทานและมีน้ำหนักเบาแต่มีราคาสูง สามารถล้างทำความสะอาด มักใช้เก็บอาหาร
      \item \textbf{พาเลทยาง} เป็นพาเลทภายในมีน้ำหนักเบาและคงทน มักใช้อุตสาหกรรมอิเล็กทรอนนิกส์เพื่อป้องกันไฟฟ้าสถิต
      \item \textbf{พาเลทกระดาษ} เป็นพาเลทขนส่งมีน้ำหนักเบาราคาถูก ใช้งานได้ 1-3 ครั้ง มักกับการขนส่งม้วนผ้า หรือเฟอร์นิเจอร์
      \item \textbf{penalty incurred:} circulation of tote at 7-11
      \item \textbf{Note: USDA:} US department of Agriculture
    \end{itemize}
  }
  \end{frame}

%============================================================================================================--
\begin{frame}
  \frametitle{Shape of Pallet}

  \begin{center}
  \begin{tabular}{ccc}
    \includegraphics[width=0.25\textwidth]{\figureLocation{03equipment/unitload/StringerPallet00.png}} &~~~&
    \includegraphics[width=0.25\textwidth]{\figureLocation{03equipment/unitload/BlockPallet00.png}} \\
          \includegraphics[width=0.25\textwidth]{\figureLocation{03equipment/unitload/BlockPallet01.png}}&~~~&
    \includegraphics[width=0.25\textwidth]{\figureLocation{03equipment/unitload/StringerPallet01.png}} \\
    \small{Block Pallet} & & \small{Stringer Pallet} \\
  \end{tabular}
  \end{center}

  \NextLine
  \textbf{Standard Pallets}
  \begin{description}
    \item[ISO pallets]  1000 mm $\times$ 1200 mm
    \item[US pallet]  40 in $\times$ 48 in  (1016 mm $\times$ 1219 mm) or 42 in $\times$ 48 in
    \item[Euro pallet]  800 mm $\times$ 600 mm \& 800 mm $\times$ 1200 mm
  \end{description}

  \note{
    \noindent\textbf{Purpose:} shape of pallet\par
    \noindent\textbf{Leading Question:} Why do pallet openings and geometry matter?\par
    \noindent\textbf{Teaching cue:} Check fork access and the support conditions for one proposed load.\par
    \begin{itemize}
      \item \textbf{stringer pallet:} '8'-shaped, expensive, durable, RT/FT only
      \item \textbf{block pallet:} '3'-shaped, cheap, both RT/FT and hand pallet
      \item \textbf{material:} circulation; durability, cost, industry
      \item \textbf{size \& form:} depend on handling
    \end{itemize}
  }
  \end{frame}


%============================================================================================================--
\begin{frame}
  \frametitle{Other pallets}

  \begin{center}
      \only<1|handout:1>{\includegraphics[width=0.75\textwidth]{\figureLocation{03equipment/unitload/otherPallets.jpg}}}
      \only<2|handout:0>{\includegraphics[width=0.75\textwidth]{\figureLocation{03equipment/unitload/palletStyle.png}}}
  \end{center}


  \note{
    \noindent\textbf{Purpose:} pallet has many shape/ size depending of object\par
    \noindent\textbf{Leading Question:} When would a nonstandard pallet be justified?\par
    \noindent\textbf{Teaching cue:} Identify the product or process constraint and check compatibility across handoffs.\par
    \includegraphics[width=0.50\textwidth]{\figureLocation{03equipment/unitload/palletStyle.png}}\\
    \textbf{Interesting Figures:} \\
    \begin{itemize}
      \item \textbf{Wired Cage:} ใช้สำหรับสินค้าชิ้นเล็ก สามารถทับซ้อนได้ ไม่ต้องการนับ เช่น รองเท้าคืน
      \item \textbf{Rolled Cage:} ใช้สำหรับขนถ่ายในพื้นที่เล็ก เช่น อาหาร
      \item \textbf{Drum:} ใช้สำหรับของเหลวหรือ สารเคมี
    \end{itemize}
  }
  \end{frame}

%============================================================================================================--
\begin{frame}
  \frametitle{Other Industrial Package}

  \begin{center}
    \begin{tabular}{ccc}
      \includegraphics[width=0.25\textwidth]{\figureLocation{03equipment/unitload/multiTote.png}} &
      \includegraphics[width=0.25\textwidth]{\figureLocation{03equipment/unitload/collapseTote.jpg}} &
      \includegraphics[width=0.25\textwidth]{\figureLocation{03equipment/unitload/metalBin.jpg}} \\
      \includegraphics[width=0.25\textwidth]{\figureLocation{03equipment/unitload/tray.png}} &
      \includegraphics[width=0.25\textwidth]{\figureLocation{03equipment/unitload/basket.png}} &
      \includegraphics[width=0.25\textwidth]{\figureLocation{03equipment/unitload/metalCage.jpg}} \\
    \end{tabular}
  \end{center}


  \note{
    \noindent\textbf{Purpose:} Connect Other Industrial Package to a decision and the evidence needed to explain it.\par
    \noindent\textbf{Leading Question:} What must industrial packaging protect during storage and movement?\par
    \noindent\textbf{Teaching cue:} Match damage risks to containment, stacking, and handling needs.\par
    Story about tote
    \begin{itemize}
      \item \textbf{tote:} 7-11 foldable tote with anti freeziing
      \item \textbf{air cargo:} import custom-made (fit to aircraft); AKH = module; LH = airline
      \item \textbf{container:} with A/C; isolated
    \end{itemize}
  }
  \end{frame}


%============================================================================================================--
\begin{frame}
  \frametitle{Storage \& retrieval equipments}

  \begin{center}
    \includegraphics[width=0.45\linewidth]{\figureLocation{03equipment/storage/DriveInRack.jpg}}
   \end{center}

  \begin{itemize}
    \item \textbf{Idea:} cubic space saving \& efficient retrieving
    \item \textbf{Where:} storage \& picking area
    \item \textbf{Issue:} standardization, FIFO, safety, ergonomic
    \item \textbf{Example:} floor stack, selective rack (single deep \& double deep), carrousel
  \end{itemize}

  \note{
    \noindent\textbf{Purpose:} Connect Storage \& retrieval equipments to a decision and the evidence needed to explain it.\par
    \noindent\textbf{Leading Question:} Which equipment choice changes access to a specific inventory unit?\par
    \noindent\textbf{Teaching cue:} Compare placement density with selectivity and replenishment effort.\par
    \begin{itemize}
      \item equipments some time yield less cubic space
      \item safety, damage (high price/fragil)
      \item allow to use high efficiently
    \end{itemize}
  }
  \end{frame}

%============================================================================================================--
\begin{frame}
  \frametitle{Floor stack: No equipment}

  \begin{tabular}{cc}
    \begin{minipage}{0.45\linewidth}
      \includegraphics[width=0.95\textwidth]{\figureLocation{03equipment/storage/floorStorage01.jpg}}
    \end{minipage}
    &
    \begin{minipage}{0.45\linewidth}
      \includegraphics[width=0.95\textwidth]{\figureLocation{03equipment/storage/floorStorage02.jpg}}
    \end{minipage}
  \end{tabular}


  %\NextHalfLine
  \begin{itemize}
    \item<1-|handout:1> \textbf{Idea:} stack pallets up-height
    \item<1-|handout:1> \textbf{Pro:} zero investment, multiple pallets per SKU, high inventory over \NextHalfLine
    \item<1-|handout:1> \textbf{Con:} \alert{honeycombing problem}, stability
    \item<1-|handout:1> \textbf{Issue:} stack-ability, stack height, aisle width
  \end{itemize}

  \note{
    \noindent\textbf{Purpose:} Connect Floor stack: No equipment to a decision and the evidence needed to explain it.\par
    \noindent\textbf{Leading Question:} When does floor stacking become impractical?\par
    \noindent\textbf{Teaching cue:} Consider stack stability, selectivity, and the sequence for accessing stock.\par
    \begin{itemize}
      \item block stacking can result in the minimum cost of storage b/c high cube utilization \& no storage medium
      \item material accessibility is low b/c only the top of the front stack is accessible
      \item \textbf{model:} trade off between accessibility VS space utilization similar to EOQ
    \end{itemize}
  }
  \end{frame}


%============================================================================================================--
\begin{frame}
  \frametitle{Single-deep rack}

  \begin{tabular}{cc}
    \begin{minipage}{0.45\linewidth}
      \includegraphics[width=0.95\textwidth]{\figureLocation{03equipment/storage/singleDeepRack3D.jpg}}
    \end{minipage}
    &
    \begin{minipage}{0.45\linewidth}
      \includegraphics[width=0.95\textwidth]{\figureLocation{03equipment/storage/singleDeepRack.jpg}}
    \end{minipage}
  \end{tabular}


  \begin{itemize}
    \item<1-|handout:1> \textbf{Idea:} a pallet rack that has a \textbf{single} storage space
    \item<1-|handout:1> \textbf{Pro:} each pallet is independently accessible
    \item<1-|handout:1> \textbf{Con:} too many aisles $\rightarrow$ inefficient space utilization
  \end{itemize}


  \note{
    \noindent\textbf{Purpose:} Connect Single-deep rack to a decision and the evidence needed to explain it.\par
    \noindent\textbf{Leading Question:} What do we gain by using single-deep rack?\par
    \noindent\textbf{Teaching cue:} Compare direct access and space use against deeper arrangements.\par
    \begin{itemize}
      \item Most popular type of storage rack
      \item Pallet is between load-supporting beams
      \item special attachments can be used to make the racks capable of supporting coils, drums, skids
    \end{itemize}
  }
  \end{frame}

%============================================================================================================--
\begin{frame}
  \frametitle{Double-deep rack}

  \begin{center}
  \begin{tabular}{cc}
    \begin{minipage}{0.45\linewidth}
      \includegraphics[width=0.95\textwidth]{\figureLocation{03equipment/storage/doubleDeepRack3D.jpg}}
    \end{minipage}
    &
    \begin{minipage}{0.45\linewidth}
      \includegraphics[width=0.95\textwidth]{\figureLocation{03equipment/storage/doubleDeepRack.jpg}}
    \end{minipage}
  \end{tabular}
  \end{center}

  \begin{itemize}
    \item \textbf{Idea:} a pallet rack that has a \textbf{double} storage space
    \item \textbf{Important:} Each lane dedicated to one SKU (one pallet or two pallets)
    \item \textbf{Pro:} Less aisle space required (upto 50\% savings in aisle space)
    \item \textbf{Con:} More work and/or specialized equipment for retrieving
  \end{itemize}

  \note{
    \noindent\textbf{Purpose:} Connect Double-deep rack to a decision and the evidence needed to explain it.\par
    \noindent\textbf{Leading Question:} What stock profile can work with double-deep rack?\par
    \noindent\textbf{Teaching cue:} Identify the access constraint and contrast dense stock with high assortment.\par
    \begin{itemize}
      \item $<<$empty$>>$
    \end{itemize}
  }
  \end{frame}



%============================================================================================================--
\begin{frame}
  \frametitle{Pallet flow rack}

  \begin{center}
  \begin{tabular}{cc}
    \begin{minipage}{0.45\linewidth}
      \includegraphics[width=0.75\textwidth]{\figureLocation{03equipment/storage/palletFlowRack3D.jpg}}
    \end{minipage}
    &
    \begin{minipage}{0.45\linewidth}
      \includegraphics[width=0.75\textwidth]{\figureLocation{03equipment/storage/palletFlowRack.jpg}}
    \end{minipage}
  \end{tabular}
  \end{center}

  \NextHalfLine
  \begin{itemize}
    \item<1-|handout:1> \textbf{Idea:} a pallet rack that \textbf{always} brings next pallet
    \item<1-|handout:1> \textbf{Important:} separate picking \& put-away
    \item<1-|handout:1> \textbf{Pro:} high pick density, FIFO
    \item<1-|handout:1> \textbf{Con:} space utilization, high cost
  \end{itemize}

  \note{
    \noindent\textbf{Purpose:} Connect Pallet flow rack to a decision and the evidence needed to explain it.\par
    \noindent\textbf{Leading Question:} What stock rotation can properly configured pallet flow support?\par
    \noindent\textbf{Teaching cue:} Trace the replenishment side, picking side, and the oldest load.\par
    \begin{itemize}
      \item typical found at very large DC
    \end{itemize}
  }
  \end{frame}

%============================================================================================================--
\begin{frame}
  \frametitle{Gravity flow rack}

  \begin{center}
  \begin{tabular}{cc}
    \begin{minipage}{0.45\linewidth}
      \includegraphics[width=0.95\textwidth]{\figureLocation{03equipment/storage/flowRack01.jpg}}
    \end{minipage}
    &
    \begin{minipage}{0.45\linewidth}
      \includegraphics[width=0.95\textwidth]{\figureLocation{03equipment/storage/flowRack02.jpg}}
    \end{minipage}
  \end{tabular}
  \end{center}

  \NextHalfLine
  \begin{itemize}
    \item<1-|handout:1> \textbf{Idea:} a rack that \textbf{always} brings next case/carton (200+ picks/hr)
    \item<1-|handout:1> \textbf{Pro:} high pick density, FIFO
    \item<1-|handout:1> \textbf{Con:} space utilization, high cost
  \end{itemize}

  \note{
    \noindent\textbf{Purpose:} Connect Gravity flow rack to a decision and the evidence needed to explain it.\par
    \noindent\textbf{Leading Question:} What handling-unit or load constraints apply to gravity flow?\par
    \noindent\textbf{Teaching cue:} Match the configuration to the original equipment documentation and safe use limits.\par
    \begin{itemize}
      \item a.k.a carton flow rack
      \item<1-|handout:1> \textbf{Where:} date-sensitive facility, fast pick area
      \item \textbf{Note:} found in DKSH rearrange every 3 months
    \end{itemize}
  }
  \end{frame}

%============================================================================================================--
\begin{frame}
  \frametitle{Bin shelve}

  \begin{center}
  \begin{tabular}{cc}
    \begin{minipage}{0.45\linewidth}
      \includegraphics[width=0.95\textwidth]{\figureLocation{03equipment/storage/binShelving.jpg}}
    \end{minipage}
    &
    \begin{minipage}{0.45\linewidth}
      \includegraphics[width=0.75\textwidth]{\figureLocation{03equipment/storage/highbayShelving.jpg}}
    \end{minipage}
  \end{tabular}
  \end{center}

  \NextHalfLine
  \begin{itemize}
    \item<1-|handout:1> \textbf{What:} storing cabinet for case/carton
    \item<1-|handout:1> \textbf{Pro:} cheap,
    \item<1-|handout:1> \textbf{Con:} single access, ID, low pick density, LIFO
  \end{itemize}

  \note{
    \noindent\textbf{Purpose:} Connect Bin shelve to a decision and the evidence needed to explain it.\par
    \noindent\textbf{Leading Question:} How does bin shelving trade accessibility against density?\par
    \noindent\textbf{Teaching cue:} Compare reach, slot size, and replenishment frequency for small items.\par
    \begin{itemize}
      \item \textbf{benefit:} use hight, reduce damage, stock on top
      \item \textbf{golden zone:}
    \end{itemize}
  }
  \end{frame}

%============================================================================================================--
\begin{frame}
  \frametitle{Small storage item equipments}

  \begin{center}
  \begin{tabular}{cc}
    \begin{minipage}{0.45\linewidth}
      \includegraphics[width=0.85\textwidth]{\figureLocation{03equipment/storage/modularDrawer.jpg}}
    \end{minipage}
    &
    \begin{minipage}{0.45\linewidth}
      \includegraphics[width=0.75\textwidth]{\figureLocation{03equipment/storage/hiDensityBin.jpg}}
    \end{minipage}
  \end{tabular}
  \end{center}

  \note{
    \noindent\textbf{Purpose:} Connect Small storage item equipments to a decision and the evidence needed to explain it.\par
    \noindent\textbf{Leading Question:} Which features make small items easier to identify and replenish?\par
    \noindent\textbf{Teaching cue:} Compare labels, slot dimensions, and pick confirmation for two options.\par
    \begin{itemize}
      \item $<<$empty$>>$
    \end{itemize}
  }
  \end{frame}

%============================================================================================================--
\begin{frame}
  \frametitle{Other Automatic Storage Equipment}

  \begin{center}
    \begin{tabular}{ccc}
      \includegraphics[width=0.25\textwidth]{\figureLocation{03equipment/storage/unitload.png}} &
      \includegraphics[width=0.25\textwidth]{\figureLocation{03equipment/storage/miniload.png}} &
      \includegraphics[width=0.25\textwidth]{\figureLocation{03equipment/storage/vShutter.png}} \\
      \small{Unitload ASRS} & \small{Miniload ASRS} & \small{Vertical Shutter} \\
    \end{tabular}
  \end{center}

  \NextHalfLine
  \begin{itemize}
    \item \textbf{Idea:} combine put-away, picking, storing
    \item \textbf{Pro:} high hight, little labor
    \item \textbf{Con:} investment, may double handling
  \end{itemize}

  \note{
    \noindent\textbf{Purpose:} introduce auto. Storage \& MHE\par
    \noindent\textbf{Leading Question:} Which operating conditions might support automated storage?\par
    \noindent\textbf{Teaching cue:} Check rate, variability, dimensions, downtime response, and cost.\par
    \textbf{Trivial:}       \\
    \begin{itemize}
      \item \textbf{Operation:} miniload (pallet) \& unitload (tote) uses stacker cranes, but vert. shutter uses vert.lift and vehicle.
      \item \textbf{Special:} STV, pick-up-delivery, station, selective rack (unitload \& vert.shuttle)
      \item \textbf{Note:} much more MHE
    \end{itemize}
  }
  \end{frame}

%============================================================================================================--
\begin{frame}
  \frametitle{Material Handling Equipment}

  \begin{center}
    \includegraphics[width=0.50\linewidth]{\figureLocation{03equipment/transport/handTruck.jpg}}
   \end{center}

  \begin{itemize}
    \item \textbf{Idea:} moving items/SKUs
    \item \textbf{Where:} everywhere
    \item \textbf{Issue:} reach, automation, space footprint, congestion
    \item \textbf{Example:} hand truck, forklift, conveyor
  \end{itemize}

  \note{
    \noindent\textbf{Purpose:} Connect Material Handling Equipment to a decision and the evidence needed to explain it.\par
    \noindent\textbf{Leading Question:} What task is the material-handling equipment meant to improve?\par
    \noindent\textbf{Teaching cue:} Name movement, lift, frequency, and load constraints before comparing equipment.\par
    \begin{itemize}
      \item \textbf{Pallet/Non-Pallet:} have forks for handling pallets or flat surface ?
      \item \textbf{Manual/Powered:} truck have manual or powered vertical (lifting) and/or horizontal (travel) movement capabilities. Manual => walk => operator provides the force needed for lifting loads and/or pushing the vehicle.
      \item \textbf{Walk/Ride:} power with walking speed
      \item \textbf{Stack/No Stack:} truck be used to lift loads for stacking purposes.
    \end{itemize}
  }
  \end{frame}

%============================================================================================================--
\begin{frame}
  \frametitle{Counter balance forklift truck}

  \begin{tabular}{cc}
    \begin{minipage}{0.45\linewidth}
      \includegraphics[width=0.95\textwidth]{\figureLocation{03equipment/transport/typicalForkLift02.jpg}}
    \end{minipage}
    &
    \begin{minipage}{0.45\linewidth}
      \includegraphics[width=0.95\textwidth]{\figureLocation{03equipment/transport/forkLift.jpg}}
    \end{minipage}
  \end{tabular}

  \NextHalfLine
  \begin{itemize}
    \item<1-|handout:1> \textbf{Idea:} unit-load mover equipped with \alert{motor} \& \alert{hydriodic}
    \item<1-|handout:1> \textbf{Pro:} very useful
    \item<1-|handout:1> \textbf{Con:} wide turn $\rightarrow$ wide aisle
  \end{itemize}

  \note{
    \noindent\textbf{Purpose:} Connect Counter balance forklift truck to a decision and the evidence needed to explain it.\par
    \noindent\textbf{Leading Question:} Which specifications must be verified before selecting a forklift?\par
    \noindent\textbf{Teaching cue:} Check rated load, center of gravity, aisle geometry, surface conditions, and training requirements.\par
    \begin{itemize}
      \item \textbf{width} 20 inches or 27 inches
      \item \textbf{weight} 1-2 tons (depend height of pallet)
      \item \textbf{left height} 2 meter
      \item \textbf{operating clearance:} 15 centimeter
      \item<1-|handout:1> \textbf{Components:} mast, fork, driving
      \item \textbf{Note:} UF used forklift as hammer
    \end{itemize}
  }
  \end{frame}

%============================================================================================================--
\begin{frame}
  \frametitle{Manual Equipment}

  \begin{center}
    \begin{tabular}{ccc}
      \includegraphics[width=0.25\textwidth]{\figureLocation{03equipment/transport/dolly.png}} &
      \includegraphics[width=0.25\textwidth]{\figureLocation{03equipment/transport/handjack.png}} &
      \includegraphics[width=0.25\textwidth]{\figureLocation{03equipment/transport/liftCart.png}} \\
    \end{tabular}
  \end{center}

  \NextHalfLine
  \begin{itemize}
    \item \textbf{Idea:} manual equipment for moving pallet or tote (no driving cab)
    \item \textbf{Pro:} small, cheap
    \item \textbf{Con:} more manual, fixed height (i.e., not apply for reversible pallet)
  \end{itemize}

  \note{
    \noindent\textbf{Purpose:} Connect Manual Equipment to a decision and the evidence needed to explain it.\par
    \noindent\textbf{Leading Question:} Which tasks justify manual equipment instead of a powered unit?\par
    \noindent\textbf{Teaching cue:} Compare loads, movement frequency, aisle access, and ergonomics.\par
    \begin{itemize}
      \item Front wheels are mounted inside the end of the forks
      \item extend to the floor as the pallet is only lifted enough to clear the floor for subsequent travel
      \item Pallet restrictions: reversible pallets cannot be used
    \end{itemize}
  }
  \end{frame}

%============================================================================================================--
\begin{frame}
  \frametitle{Industrial Trucks}

  \begin{center}
    \begin{tabular}{ccc}
      \includegraphics[width=0.25\textwidth]{\figureLocation{03equipment/transport/swingMast.png}} &
      \includegraphics[width=0.25\textwidth]{\figureLocation{03equipment/transport/rt.png}} &
      \includegraphics[width=0.25\textwidth]{\figureLocation{03equipment/transport/vna.png}} \\
      \small{Swing Mast} & \small{Reach Truck} & \small{VNA truck} \\
    \end{tabular}
  \end{center}

  \NextHalfLine
  \begin{itemize}
    \item \textbf{Idea:} moving pallet from A $\rightarrow$ B with power
    \item \textbf{Type:} turret, footprint, drivable, \# pallets
    \item \textbf{Pro:} save time \& labor
    \item \textbf{Con:} price, storage equipment
  \end{itemize}

  \note{
    \noindent\textbf{Purpose:} Connect Industrial Trucks to a decision and the evidence needed to explain it.\par
    \noindent\textbf{Leading Question:} How can aisle geometry and lift requirements exclude an industrial truck?\par
    \noindent\textbf{Teaching cue:} Match the equipment specification to available clearance and load.\par
    \begin{itemize}
      \item 8-10 ft. minimum aisle width requirement
      \item extended reaching mechanisms are available to enable double-deep storage
      \item electric power; camera; balance \& stability
    \end{itemize}
  }
  \end{frame}


%============================================================================================================--
\begin{frame}
  \frametitle{Crane}

  \begin{tabular}{cc}
    \begin{minipage}{0.45\linewidth}
      \includegraphics[width=0.95\textwidth]{\figureLocation{03equipment/transport/craneDrawing.jpg}}
    \end{minipage}
    &
    \begin{minipage}{0.45\linewidth}
      \includegraphics[width=0.95\textwidth]{\figureLocation{03equipment/transport/asrs.png}}
    \end{minipage} \\
    \tiny{bridge crane} & \tiny{stacker crane in AS/RS}
  \end{tabular}

  \NextHalfLine
  \begin{itemize}
    \item \textbf{Idea:} moving items \alert{overhead}
    \item \textbf{Pro:} flexible shape/size
    \item \textbf{Con:} restricted area, congestion with others
  \end{itemize}

  \note{
    \noindent\textbf{Purpose:} introduction to crane\par
    \noindent\textbf{Leading Question:} What workplace constraint affects crane selection?\par
    \noindent\textbf{Teaching cue:} Identify load limits, travel path, lifting frequency, and site requirements.\par
    \textbf{Interesting Figures \& Trivial:} \\
    \begin{itemize}
      \item \textbf{เครนแบบสะพาน} (Bridge Crane) เป็นเครนที่เคลื่อนที่ในแนวดิ่งด้วย (Hoist) ซึ่งเป็นระบบรอกพร้อมขอเกี่ยววัสดุ
      \item \textbf{เครนแบบยกซ้อน} (Stacker Crane) เป็นแครนพร้อมแท่นแพลตฟอร์ม (Platform) เคลื่อนที่ในแนวราบโดยรางโลหะซึ่งอาจจะติดตั้งระดับพื้นอาคารและ/หรือเพดาน การเคลื่อนที่ในแนวดิ่งโดยการปรับตำแหน่งแทนแพลตฟอร์มด้วยมอเตอร์ไฟฟ้า  เครนแบบยกซ้อนมักประยุกต์ใช้เป็นส่วนหนึ่งของระบบ AS/RS
      \item \textbf{Operation:} move loads horizontally \& vertically over a restricted area
      \item \textbf{Justification} for insufficient (or intermittent) volume, not good for conveyor
      \item \textbf{PATH:} conveyor $<$ crane $<$ truck
    \end{itemize}
  }
  \end{frame}



%============================================================================================================--
\begin{frame}
  \frametitle{Others MHE: Rail}

  \begin{center}
  \begin{tabular}{ccc}
    \includegraphics[width=0.25\textwidth]{\figureLocation{03equipment/transport/towConveyor.png}} &
    \includegraphics[width=0.25\textwidth]{\figureLocation{03equipment/transport/agv.png}} &
    \includegraphics[width=0.25\textwidth]{\figureLocation{03equipment/transport/rgv.png}} \\
      \tiny{Tow Line} & \tiny{AGV} & \tiny{RGV} \\
  \end{tabular}
  \end{center}

  \NextHalfLine
  \begin{itemize}
    \item \textbf{Idea:} moving items on \alert{fixed paths} usually as loop
    \item \textbf{Pro:} eliminate worker
    \item \textbf{Con:} restricted area, investment
  \end{itemize}


  \note{
    \noindent\textbf{Purpose:} special MHE\par
    \noindent\textbf{Leading Question:} When could a fixed rail path improve material movement?\par
    \noindent\textbf{Teaching cue:} Compare repeatable flow and route volume with the need for flexible routing.\par
   \includegraphics[width=0.45\textwidth]{\figureLocation{03equipment/transport/railMHE.png}}
   \includegraphics[width=0.45\textwidth]{\figureLocation{03equipment/transport/conveyorMHE.png}} \\
    \textbf{Interesting Figures \& Trivial:} \\
    \begin{itemize}
      \item \textbf{Tow Conveyor} instead of just the forks, the entire mast rotates
      \item \textbf{AGV} auto.guide vehicle/ not restricted, buy easy to obstruct/ cheaper but slow
      \item \textbf{RGV} rail.guide vehicle/ restricted to path/ more expensive $\approx$ 1.0 M THB / unit

    \end{itemize}
  }
  \end{frame}

%============================================================================================================--
\begin{frame}
  \frametitle{Conveyor}

  \begin{center}
  \begin{tabular}{cc}
    \begin{minipage}{0.45\linewidth}
      \includegraphics[width=0.95\textwidth]{\figureLocation{03equipment/transport/conveyor01.jpg}}
    \end{minipage}
    &
    \begin{minipage}{0.45\linewidth}
      \includegraphics[width=0.95\textwidth]{\figureLocation{03equipment/transport/conveyor02.jpg}}
    \end{minipage}
  \end{tabular}
  \end{center}

  \NextHalfLine
  \begin{itemize}
    \item<1-|handout:1> \textbf{What:} automatic moving 'regular' shape pallet
    \item<1-|handout:1> \textbf{Pro:} free labor
    \item<1-|handout:1> \textbf{Con:} large moving huge std. size, fixed paths
  \end{itemize}

  \note{
    \noindent\textbf{Purpose:} Connect Conveyor to a decision and the evidence needed to explain it.\par
    \noindent\textbf{Leading Question:} Which flow conditions make a conveyor worthwhile?\par
    \noindent\textbf{Teaching cue:} Check continuity, volume, handoffs, and peak congestion before selecting it.\par
    \begin{itemize}
      \item<1-|handout:1> \textbf{Example:} chute, roll, belt, sortation conveyor
      \item \textbf{load:} unit/bulk
      \item \textbf{location:} on floor/in floor/overhead
      \item \textbf{wait:} accumulate/ no accumulate
    \end{itemize}
    Goal: move a lot; between known \& fixed locations
    \begin{itemize}
      \item \textbf{chute conveyor:} used to link two handling devices (accumulated)
      \item \textbf{roll conveyor:} powered or non-powered (accumulate)
      \item \textbf{belt conveyor:} transporting light- \& medium-weight loads between operations (no accumulate)
      \item \textbf{sortition conveyor:} system of belt convey
    \end{itemize}
  }
  \end{frame}



%============================================================================================================--
\begin{frame}
  \frametitle{Variation of Conveyors}


  \begin{center}
    \begin{tabular}{ccc}
      \includegraphics[width=0.25\textwidth]{\figureLocation{03equipment/transport/beltConveyor.png}} &
      \includegraphics[width=0.20\textwidth]{\figureLocation{03equipment/transport/chainConveyor.png}} &
      \includegraphics[width=0.25\textwidth]{\figureLocation{03equipment/transport/gravConveyor.png}} \\[1pt]
       \tiny{belt} & \tiny{chain} & \tiny{gravity} \\[1pt]
      \multicolumn{2}{c}{\includegraphics[width=0.50\textwidth]{\figureLocation{03equipment/transport/screwConveyor.png}}} &
      \includegraphics[width=0.25\textwidth]{\figureLocation{03equipment/transport/troughedConveyor.png}} \\[1pt]
       \tiny{screw} &  & \tiny{troughted } \\
    \end{tabular}
  \end{center}

  \begin{itemize}
    \item \textbf{Selection:} materials, slope, price, weight
  \end{itemize}

  \note{
    \noindent\textbf{Purpose:} intro to conveyor\par
    \noindent\textbf{Leading Question:} Why might different conveyor types suit different packages?\par
    \noindent\textbf{Teaching cue:} Match the contact surface and package geometry to the product and routing needs.\par
    \textbf{Interesting Figures \& Trivial:} \\
    \begin{itemize}
      \item \textbf{belt conveyor:} transporting light \& medium-weight loads between operations (no accumulate), move up or down
      \item \textbf{chain conveyor:} use for transfer pallets
      \item \textbf{gravity conveyor:} use for roller conveyor
      \item \textbf{screw conveyor:} used for transfer bulk items
      \item \textbf{troughes conveyor:} use for transfer bulk items, 30degree
      \item \textbf{chute conveyor:} use for link two handling devices (accumulated)
      \item \textbf{roll conveyor:} powered or non-powered (accumulate)
      \item \textbf{sortation conveyor:} system of belt convey
    \end{itemize}
  }
  \end{frame}

%============================================================================================================--
\begin{frame}
  \frametitle{Identification equipment}

  \begin{center}
    \includegraphics[width=0.50\linewidth]{\figureLocation{03equipment/it/rfid.jpg}}
   \end{center}

  \begin{itemize}
    \item \textbf{Idea:} speeding receiving \& shipping
    \item \textbf{Where:} receiving \& shipping
    \item \textbf{Issue:} integration with system
    \item \textbf{Example:} RFID, bar code reader, magnet
  \end{itemize}

  \note{
    \noindent\textbf{Purpose:} Connect Identification equipment to a decision and the evidence needed to explain it.\par
    \noindent\textbf{Leading Question:} Which events require a reliable item or location identifier?\par
    \noindent\textbf{Teaching cue:} Trace the identifier through receiving, storage, picking, and dispatch.\par
    \begin{itemize}
      \item \textbf{Magnet:} Data encoded on a magnetic stripe that is readable in almost any environment; Requires contact with a reader
      \item \textbf{RF-ID:} Noncontact: can be read when the tag is within 30 ft. of an antenna
    \end{itemize}
  }
  \end{frame}

%============================================================================================================--
\begin{frame}
  \frametitle{Useful information }

  \textbf{Material Handling}
    \begin{itemize}
      \item \textbf{Material Handling Taxonomy:} \href{http://www.mhia.org/industrygroups/cicmhe/resources/mhe_tax.htm}{\code{http://www.mhia.org/industrygroups/cicmhe/resources/mhe\_tax.htm}}
      \item \textbf{Material Handling Pictures:} \href{https://www.cirrelt.ca/mhmultimediabank/}{\code{https://www.cirrelt.ca/mhmultimediabank/}}
      %\item \textbf{Material Handling Guide:} Chula Blackboard
    \end{itemize}

  \NextLine
  \textbf{Warehouse Tours}
  \begin{itemize}
    \item \textbf{Interactive Tour:} \href{http://www.roodbergen.com/warehouse/}{\code{http://www.roodbergen.com/warehouse/}}
    \item \textbf{Warehouse Science:} \href{http://www2.isye.gatech.edu/~jjb/wh/sites/sites.html}{\code{http://www2.isye.gatech.edu/\~jjb/wh/sites/sites.html}}
    %\item \textbf{Multimedia Tour:} Chula Blackboard
  \end{itemize}


  \note{
    \noindent\textbf{Purpose:} Connect Useful information to a decision and the evidence needed to explain it.\par
    \noindent\textbf{Leading Question:} Which information would an operator need at the point of work?\par
    \noindent\textbf{Teaching cue:} Name item, location, quantity, work priority, and what to do with an exception.\par
    \begin{itemize}
      \item $<<$empty$>>$
    \end{itemize}
  }
  \end{frame}

%------------------------------------------------------
\tikzset{MyArrow/.style={single arrow, draw, minimum width=15pt, minimum height=30pt,
                         inner sep=3pt, single arrow head extend=5pt,fill=black!30}
}
\tikzstyle{opBlock}=[text centered,text height=1.0mm,opacity=.5,text opacity=1,minimum width=45pt]
\tikzstyle{path}=[rectangle]
\tikzstyle{demand}=[fill=black!35,pattern=north west lines]
\tikzstyle{demandRed}=[fill=red!90]
\tikzstyle{demandBlu}=[fill=blue!90]
\usetikzlibrary{patterns}
\tikzset{diagonal fill/.style 2 args={fill=#2, path picture={
\fill[#1, sharp corners] (path picture bounding box.south west) -|
                         (path picture bounding box.north east) -- cycle;}},
reversed diagonal fill/.style 2 args={fill=#2, path picture={
\fill[#1, sharp corners] (path picture bounding box.north west) |-
                         (path picture bounding box.south east) -- cycle;}}
}
%%============================================================================================================--
\section[Layout]{Typical Warehouse Layout}
%%============================================================================================================--

\begin{frame}
  \frametitle{U-Shaped Layout}


  \begin{center}
     \begin{tikzpicture}[scale=.22]
     \tikzstyle{dock}=[fill=black!15];

      \draw[black,thick] (-2.5,-.5) -- (-2.5,18.5) -- (28.5,18.5) -- (28.5,-.5) -- cycle ;

      \fill[fill=green!90] (-2, 0) rectangle ++(1,2);
      \fill[demand] (-2, 2) rectangle ++(1,8);
      \fill[demandRed] (-2,10) rectangle ++(1,2);
      \draw (-2, 0) grid (-1,12);

      \foreach \x in {0,...,8}
        {
          \fill[fill=green!90] ( 3*\x, 0) rectangle ++(2,2);
          \fill[demand] ( 3*\x, 2) rectangle ++(2,8);
          \fill[demandRed] ( 3*\x,10) rectangle ++(2,2);
          \draw (3*\x, 0) grid (3*\x+2,12);
        }

      \fill[fill=green!90] (27, 0) rectangle ++(1,2);
      \fill[demand] (27, 2) rectangle ++(1,8);
      \fill[demandRed] (27,10) rectangle ++(1,2);
      \draw (27, 0) grid (28,12);

      \node[] (shipping) at ( 4,16) {\textbf{\small{Shipping Area}}};
      \node[] (reciving) at (22,16) {\textbf{\small{Receiving Area}}};

      \node[opBlock,fill=white!30] (fastItem) at (13,11) {\scriptsize{Fast-Moving}};
      \node[opBlock,fill=white!30] (mediItem) at (13, 6) {\scriptsize{Median-Moving}};
      \node[opBlock,fill=white!30] (slowItem) at (13, 1) {\scriptsize{Slow-Moving}};


      \foreach \x in {1,...,3}
      {
        \path (-1+3*\x,22.5) ++(0,-2) node[rotate= 90,MyArrow]  {};
        \path (16+3*\x,22.5) ++(0,-2) node[rotate=-90,MyArrow]  {};
        \fill[dock] (-5+4*\x,18.35) rectangle ++(2.5,0.3);
        \fill[dock] ( 7+4*\x,18.35) rectangle ++(2.5,0.3);
      }
      \fill[dock] ( 23,18.35) rectangle ++(2.5,0.3);
    \end{tikzpicture}
  \end{center}

  \begin{itemize}
      \item \textbf{Idea:} single I/O point, very convenient locations
      \item \textbf{Where:} distribution network
      \item \textbf{Issue:} congestion near receiving \& shipping
  \end{itemize}

  \note{
    \noindent\textbf{Purpose:} Connect U-Shaped Layout to a decision and the evidence needed to explain it.\par
    \noindent\textbf{Leading Question:} What makes the U-layout suitable and where can traffic conflict occur?\par
    \noindent\textbf{Teaching cue:} Trace receiving and shipping flows at a busy hour and identify one mitigation.\par
    \begin{itemize}
      \item \textbf{Overview:} layout can tell you type of WH, unit it handles
      \item \textbf{CAT by flow:} 'U' \& 'I' (if only one storage area)
      \item \textbf{Very convenient:} near I/O travels small; distant from I/O travels very far
      \item \textbf{Benefits:} share admin; sharing work receiving \& shipping; combine works; control
      \item \textbf{Good for DC:} DC (fluc $<$ scale $<$ break bulk $<$ consolidate); focus on move item fast
      \item \textbf{Idea:} single I/O point, very convenient locations
      \item \textbf{Where:} distribution network
      \item \textbf{Issue:} congestion near receiving \& shipping
    \end{itemize}
  }
  \end{frame}
%============================================================================================================--

\begin{frame}
  \frametitle{I-Shaped Layout}

  \begin{center}
     \begin{tikzpicture}[scale=.22]
     \tikzstyle{dock}=[fill=black!15];

      \draw[black,thick] (-2.5,-.5) -- (-2.5,18.5) -- (28.5,18.5) -- (28.5,-.5) -- cycle ;

      \fill[demand] ( 2, 0) rectangle ++(22,1);
      \draw ( 2, 0) grid ( 24,1);


      \foreach \x in {0,...,4}
        {
        \fill[demand] ( 2, 3*\x+2) rectangle ++(22,2);
        \draw ( 2, 3*\x+2) grid ( 24,3*\x+4);
        }

      \fill[demand] ( 2, 17) rectangle ++(22,1);
      \draw ( 2, 17) grid ( 24,18);

      \node[rotate=90] at (26, 9) {\textbf{\small{Shipping Area}}};
      \node[rotate=-90] at ( 0, 9) {\textbf{\small{Recieving Area}}};
      \node[opBlock,fill=white!30] at (13, 9) {\scriptsize{Median-Moving}};



      \foreach \x in {1,...,3}
      {
        \path (-4.0,4*\x+1) ++(-2,0) node[anchor=west,MyArrow]  {};
        \path (31.0,4*\x+1) ++(-2,0) node[anchor=west,MyArrow]  {};
        \fill[dock] (-2.65,-3+5*\x) rectangle ++(0.3,2.5);
        \fill[dock] (28.35,-3+5*\x) rectangle ++(0.3,2.5);
      }

    \end{tikzpicture}
  \end{center}

  \begin{itemize}
        \item \textbf{Idea:} two I/O point, many convenient locations
        \item \textbf{Where:} manufacturing warehouse
        \item \textbf{Issue:} traveling distance, \# of admin, \& combine works
  \end{itemize}


  \note{
    \noindent\textbf{Purpose:} Connect I-Shaped Layout to a decision and the evidence needed to explain it.\par
    \noindent\textbf{Leading Question:} When can an I-shaped layout separate conflicting traffic?\par
    \noindent\textbf{Teaching cue:} Compare directional flow and available dock or yard access with a U-layout.\par
    \begin{itemize}
      \item \textbf{Many convenient:} long aisle distances are the same
      \item \textbf{Benefits:} simple flow,
      \item \textbf{Good for Manuf:} constant prod
        \item \textbf{Idea:} two I/O point, many convenient locations
        \item \textbf{Where:} manufacturing warehouse
        \item \textbf{Issue:} traveling distance, \# of admin, \& combine works
    \end{itemize}
  }
  \end{frame}

%============================================================================================================--
\begin{frame}
  \frametitle{Custom Layout}


 \begin{center}
     \begin{tikzpicture}[scale=.22]

      \draw[black,thick] (-2.5,-.5) -- (-2.5,18.5) -- (28.5,18.5) -- (28.5,-.5) -- cycle ;
      \draw[black,thin,fill=yellow!30] (0,0.0) -- (0,16.0) -- (8.5,16.0) -- (8.5,0.0) -- cycle ;

      \fill[demand] (10, 0) rectangle ++(1,5);
      \fill[demand] (10,13) rectangle ++(1,5);
      \fill[demand] (12, 0) rectangle ++(2,8);
      \fill[demand] (12,10) rectangle ++(2,8);
      \fill[demand] (15, 0) rectangle ++(2,8);
      \fill[demand] (15,10) rectangle ++(2,8);
      \fill[demand] (18, 10) rectangle ++(2,8);
      \fill[demand] (18,  0) rectangle ++(2,8);
      \fill[demand] (21, 10) rectangle ++(1,8);
      \fill[demand] (21,  0) rectangle ++(1,8);

      \fill[demandRed] (9,0) rectangle ++(1,8);
      \fill[demandRed] (9,10) rectangle ++(1,8);
      \fill[demandRed] (10,5) rectangle ++(1,3);
      \fill[demandRed] (10,10) rectangle ++(1,3);
      \fill[demandRed] (12,5) rectangle ++(1,3);
      \fill[demandRed] (12,10) rectangle ++(1,3);
      \fill[demandRed] (13,6) rectangle ++(1,2);
      \fill[demandRed] (13,10) rectangle ++(1,2);
      \fill[demandRed] (15,6) rectangle ++(1,2);
      \fill[demandRed] (15,10) rectangle ++(1,2);



      %\fill[fill=green!90] (13,12) rectangle ++(1,6);


      \fill[fill=green!90] ( 22, 10) rectangle ++(1,8);
      \fill[fill=green!90] ( 22,  0) rectangle ++(1,8);
      \fill[fill=green!90] ( 24, 10) rectangle ++(2,8);
      \fill[fill=green!90] ( 24,  0) rectangle ++(2,8);
      \fill[fill=green!90] ( 27, 10) rectangle ++(1,8);
      \fill[fill=green!90] ( 27,  0) rectangle ++(1,8);

      \foreach \x in {3,...,8}
        {
          \draw (3*\x, 10) grid (3*\x+2,18);
          \draw (3*\x, 0) grid (3*\x+2,8);
        }
      \draw (27, 10) grid (28,18);
      \draw (27, 0) grid (28,8);

      \node[] (shipping) at ( 4,17) {\textbf{\small{Shipping Area}}};
      \node[black!70] (crossing) at ( 4,9) {\begin{tabular}{c}\textbf{\small{Crossdocking}} \\ \textbf{\small{Area}}\end{tabular}};
      \node[rotate=-90] (reciving) at (-1, 9) {\textbf{\small{Recieving Area}}};

      \node[opBlock,fill=white!30] (fastItem0) at (12, 7) {\scriptsize{Fast-Moving}};
      \node[opBlock,fill=white!30] (fastItem0) at (12, 11) {\scriptsize{Fast-Moving}};
      \node[opBlock,fill=white!30] (mediItem0) at (16, 3)  {\scriptsize{Median-Moving}};
      \node[opBlock,fill=white!30] (mediItem1) at (16, 15) {\scriptsize{Median Moving}};
      \node[opBlock,fill=white!30] at (25, 9) {\begin{tabular}{c} \scriptsize{Slow-Moving} \\ \scriptsize{or Non-Moving}\end{tabular} };


      \foreach \x in {1,...,3}
      {
        \path (-1+3*\x,22) ++(0,-2) node[rotate= 90,MyArrow]  {};
        \path (-3.5,4*\x+1) ++(-2,0) node[anchor=west,MyArrow]  {};
      }

    \end{tikzpicture}
 \end{center}

  \begin{itemize}
    \item \textbf{Idea:} mixed crossdock with storage
    \item \textbf{Where:} $\ll$depends$\gg$
    \item \textbf{Issue:} $\ll$depends$\gg$
  \end{itemize}

  \note{
    \noindent\textbf{Purpose:} Connect Custom Layout to a decision and the evidence needed to explain it.\par
    \noindent\textbf{Leading Question:} Which site constraint makes a custom layout necessary?\par
    \noindent\textbf{Teaching cue:} Sketch the site limits, movement paths, and demand profile before selecting a layout.\par
    \begin{itemize}
      \item \textbf{Overview:} majority of WH is in this cat $\rightarrow$ handling; channel; dock door
      \item \textbf{Idea:} combine `U-shaped' \& `I-shaped' (a.k.a 'L'-shaped?)
      \item \textbf{Issue:} depends, mostly traffic/congestion
    \end{itemize}
  }
  \end{frame}

%============================================================================================================--
\section[Tour]{Warehouse Intuition}
%============================================================================================================--
\begin{frame}
  \frametitle{Warehouse Tours}

    \begin{itemize}
      \item \textbf{Goals:} sharing experiences \& problems, connecting to industry
      \item \textbf{Organization:} looking pro-con, brainstorm, discussion
      \item \textbf{Observe:}
      \begin{itemize}
        \item What are units handle in each warehouse?
        \item What are equipments in each warehouse? (Is it make sense?)
        \item What are the problems/issues in each warehouse? \NextHalfLine
      \end{itemize}
      \item \textbf{Next steps:} Field trips, Project
    \end{itemize}

  \note{
    \noindent\textbf{Purpose:} Connect Warehouse Tours to a decision and the evidence needed to explain it.\par
    \noindent\textbf{Leading Question:} What can a warehouse tour show that a published KPI cannot?\par
    \noindent\textbf{Teaching cue:} Ask for an observable process detail and a metric that could validate it.\par
    \begin{itemize}
      \item organization of rack meaning
      \item aisles are not equal/ corners are more convenience
      \item \textbf{Warehouse Tour:} Amazon.com, IKEA, Liquor DC, Office Depot, WestLog DC
      \begin{itemize}
        \item How each warehouse handle items?
        \item What are equipments \& operational processes in each location? \NextHalfLine
      \end{itemize}
      \item \textbf{Equipments:} unit load, high-bay rack, carousal, Kiva, conveyor, pick-to-light, pick-to-voice
      \begin{itemize}
        \item How each equipment intergraded to the layout?
        \item What are advantages \& disadvantages of each equipments? \NextHalfLine
      \end{itemize}
      \item \textbf{Accident:} Forklift accidents
    \end{itemize}
  }

  \end{frame}
%============================================================================================================--
\begin{frame}
  \frametitle{Warehouse Intuition: Ace DC}

  \begin{center}
    \only<1|handout:1> {\includegraphics[width=0.65\linewidth]{\figureLocation{04tours/AceHardware/AceDC01.jpg}}}
  \end{center}

  \begin{center}
    \only<2|handout:0> {\includegraphics[width=0.65\linewidth]{\figureLocation{04tours/AceHardware/AceDC02.jpg}}}
    \only<3|handout:0> {\includegraphics[width=0.65\linewidth]{\figureLocation{04tours/AceHardware/AceDC03.jpg}}}
    \only<4|handout:0> {\includegraphics[width=0.65\linewidth]{\figureLocation{04tours/AceHardware/AceDC04.jpg}}}
    \only<5|handout:0> {\includegraphics[width=0.65\linewidth]{\figureLocation{04tours/AceHardware/AceDC05.jpg}}}
    \only<6|handout:0> {\includegraphics[width=0.65\linewidth]{\figureLocation{04tours/AceHardware/AceDC06.jpg}}}
    \only<7|handout:0> {\includegraphics[width=0.40\linewidth]{\figureLocation{04tours/AceHardware/AceDC07.jpg}}}
    \only<8|handout:0> {\includegraphics[width=0.65\linewidth]{\figureLocation{04tours/AceHardware/AceDC08.jpg}}}
    \only<9|handout:0> {\includegraphics[width=0.65\linewidth]{\figureLocation{04tours/AceHardware/AceDC09.jpg}}}
    \only<10|handout:0> {\includegraphics[width=0.65\linewidth]{\figureLocation{04tours/AceHardware/AceDC10.jpg}}}
    \only<11|handout:0> {\includegraphics[width=0.65\linewidth]{\figureLocation{04tours/AceHardware/AceDC11.jpg}}}
    \only<12|handout:0> {\includegraphics[width=0.65\linewidth]{\figureLocation{04tours/AceHardware/AceDC12.jpg}}}
    \only<13|handout:0> {\includegraphics[width=0.65\linewidth]{\figureLocation{04tours/AceHardware/AceDC13.jpg}}}
    \only<14|handout:0> {\includegraphics[width=0.65\linewidth]{\figureLocation{04tours/AceHardware/AceDC14.jpg}}}
    \only<15|handout:0> {\includegraphics[width=0.65\linewidth]{\figureLocation{04tours/AceHardware/AceDC15.jpg}}}
    \only<16|handout:0> {\includegraphics[width=0.65\linewidth]{\figureLocation{04tours/AceHardware/AceDC16.jpg}}}
    \only<17|handout:0> {\includegraphics[width=0.40\linewidth]{\figureLocation{04tours/AceHardware/AceDC17.jpg}}}
    \only<18|handout:0> {\includegraphics[width=0.65\linewidth]{\figureLocation{04tours/AceHardware/AceDC18.jpg}}}
    \only<19|handout:0> {\includegraphics[width=0.65\linewidth]{\figureLocation{04tours/AceHardware/AceDC19.jpg}}}
   % \only<20|handout:1> {\includegraphics[width=0.65\linewidth]{\figureLocation{AceHardware/AceDC20.jpg}}}
  \end{center}

  \NextHalfLine
  \begin{flushright}
    \tiny{\textbf{source:} Tours of warehouses, distribution centers, crossdocks. \href{http://www2.isye.gatech.edu/~jjb/wh/sites/sites.html}{\code{http://www2.isye.gatech.edu/\~jjb/wh/sites/sites.html}}\hskip10pt}
  \end{flushright}

  \note{
    \noindent\textbf{Purpose:} Connect Warehouse Intuition: Ace DC to a decision and the evidence needed to explain it.\par
    \noindent\textbf{Leading Question:} Which principle could carry from the Ace facility to a different site?\par
    \noindent\textbf{Teaching cue:} Identify a visible flow choice and the constraints that make it transferable.\par
    \small{
    \begin{itemize}
      \item Colon Free Trade Zone of the Republic of Panama. They have low labor costs but must cope with space that is surprisingly tight
      \item 25,000 skus, about 80\% arrive via freighter from China (in box carton to save space, esp October through December)
      \item expanding warehouse/with water \& congestion
    \end{itemize}
    }
  }
  \end{frame}

%============================================================================================================--
\begin{frame}
  \frametitle{Warehouse Intuition: Pepsi, Atlanta}

  \begin{center}
    \only<1|handout:1> {\includegraphics[width=0.65\linewidth]{\figureLocation{04tours/Pepsi/pepsi00.jpg}}}
  \end{center}
  \begin{center}
    \only<2|handout:0> {\includegraphics[width=0.65\linewidth]{\figureLocation{04tours/Pepsi/pepsi01.jpg}}}
    \only<3|handout:0> {\includegraphics[width=0.65\linewidth]{\figureLocation{04tours/Pepsi/pepsi02.jpg}}}
    \only<4|handout:0> {\includegraphics[width=0.65\linewidth]{\figureLocation{04tours/Pepsi/pepsi03.jpg}}}
    \only<5|handout:0> {\includegraphics[width=0.65\linewidth]{\figureLocation{04tours/Pepsi/pepsi04.jpg}}}
    \only<6|handout:0> {\includegraphics[width=0.40\linewidth]{\figureLocation{04tours/Pepsi/pepsi05.jpg}}}
    \only<7|handout:0> {\includegraphics[width=0.65\linewidth]{\figureLocation{04tours/Pepsi/pepsi06.jpg}}}
    \only<8|handout:0> {\includegraphics[width=0.65\linewidth]{\figureLocation{04tours/Pepsi/pepsi07.jpg}}}
    \only<9|handout:0> {\includegraphics[width=0.65\linewidth]{\figureLocation{04tours/Pepsi/pepsi08.jpg}}}
    \only<10|handout:0> {\includegraphics[width=0.65\linewidth]{\figureLocation{04tours/Pepsi/pepsi09.jpg}}}
    \only<11|handout:0> {\includegraphics[width=0.65\linewidth]{\figureLocation{04tours/Pepsi/pepsi10.jpg}}}
    \only<12|handout:0> {\includegraphics[width=0.65\linewidth]{\figureLocation{04tours/Pepsi/pepsi11.jpg}}}
  \end{center}

  \NextHalfLine
  \begin{flushright}
    \tiny{\textbf{source:} Tours of warehouses, distribution centers, crossdocks. \href{http://www2.isye.gatech.edu/~jjb/wh/sites/sites.html}{\code{http://www2.isye.gatech.edu/\~jjb/wh/sites/sites.html}}\hskip10pt}
  \end{flushright}

  \note{
    \noindent\textbf{Purpose:} Connect Warehouse Intuition: Pepsi, Atlanta to a decision and the evidence needed to explain it.\par
    \noindent\textbf{Leading Question:} Which aspects of the Pepsi site depend on product and shipment mix?\par
    \noindent\textbf{Teaching cue:} Compare a unit load flow with the needs of a high-assortment operation.\par
    \small{
    \begin{itemize}
      \item There are relatively few skus but they go to many destinations. Most activity is picking cartons from pallets, but without automation.
      \item This one stocks around 400 skus. Pepsi tries to keep no more than eight days of inventory on hand. The facility recently converted to “pick-to-voice” \& reports that training time has been reduced
      \item There are about 16 order pickers working within 7 aisles. The forward pick area is restocked in the morning after visual check. A popular sku will have several pallets floor-stacked
    \end{itemize}
    }
  }
  \end{frame}

 %============================================================================================================--
\begin{frame}
  \frametitle{Warehouse Intuition: THD}


  \begin{center}
    \only<1|handout:1> {
      \begin{minipage}{0.95\textwidth}
        \begin{tabular}{cc}
          \includegraphics[width=0.45\textwidth]{\figureLocation{04layout/hd/network2007.png}}
          &
          \includegraphics[width=0.45\textwidth]{\figureLocation{04layout/hd/network2014.png}}
      \end{tabular}
    \end{minipage}
    }
  \end{center}
  \begin{center}
    \only<2|handout:0> {\includegraphics[width=0.85\linewidth]{\figureLocation{04layout/hd/netork.png}}}
    \only<3|handout:0> {\includegraphics[width=0.75\linewidth]{\figureLocation{04layout/hd/layout.png}}}
    \only<4|handout:0> {\includegraphics[width=0.65\linewidth]{\figureLocation{04layout/hd/inside.png}}}

  \end{center}


  \note{
    \noindent\textbf{Purpose:} intro to Hafele (sell house related equipment, key, door knob)\par
    \noindent\textbf{Leading Question:} What does the THD example show about equipment, layout, or order flow?\par
    \noindent\textbf{Teaching cue:} Name one observation and identify the missing data needed to compare performance.\par
    \textbf{Note:} our class field trip on 2012\\
    \textbf{Delivery:} major customer furniture retailer. Index, SB\\
    \begin{itemize}
      \item \textbf{Business nature:} Hardware manufacturer/importer $\rightarrow$ Furniture
      \item \textbf{Site location:} Bang-Na, BKK
      \item \textbf{Warehouse nature:} fashionable goods; delivery
      \item \textbf{Items:} 50,000 SKUs; 80\% transported in carton from China
      \item \textbf{Note:} first cable VNA in TH, order pick Sys, very advance in term of material handling
      \item \textbf{Problem:} deadstock, fashionable goods, storage area
    \end{itemize}
    }

  \end{frame}

%============================================================================================================--
\begin{frame}
  \frametitle{Warehouse Intuition: Boon Tha Vorn}

  \begin{center}
    \only<1|handout:1> {\includegraphics[width=0.40\linewidth]{\figureLocation{04tours/BoonThaVorn/BTV_Sign.jpg}}}
  \end{center}
  \begin{center}
    \only<2|handout:0> {\includegraphics[width=0.65\linewidth]{\figureLocation{04tours/BoonThaVorn/BTV_Inside00.jpg}}}
    \only<3|handout:0> {\includegraphics[width=0.95\linewidth]{\figureLocation{04tours/BoonThaVorn/BTV_Layout01.jpg}}}
    \only<4|handout:0> {\includegraphics[width=0.65\linewidth]{\figureLocation{04tours/BoonThaVorn/BTV_Inside02.jpg}}}
    \only<5|handout:0> {\includegraphics[width=0.65\linewidth]{\figureLocation{04tours/BoonThaVorn/BTV_Inside03.jpg}}}
    \only<6|handout:0> {\includegraphics[width=0.65\linewidth]{\figureLocation{04tours/BoonThaVorn/BTV_Inside04.jpg}}}
    \only<7|handout:0> {\includegraphics[width=0.65\linewidth]{\figureLocation{04tours/BoonThaVorn/BTV_Product01.jpg}}}
    \only<8|handout:0> {\includegraphics[width=0.40\linewidth]{\figureLocation{04tours/BoonThaVorn/BTV_Inside05.jpg}}}
    \only<9|handout:0> {\includegraphics[width=0.40\linewidth]{\figureLocation{04tours/BoonThaVorn/BTV_Inside06.jpg}}}
    \only<10|handout:0> {\includegraphics[width=0.40\linewidth]{\figureLocation{04tours/BoonThaVorn/BTV_Inside07.jpg}}}
    \only<11|handout:0> {\includegraphics[width=0.40\linewidth]{\figureLocation{04tours/BoonThaVorn/BTV_Inside08.jpg}}}
    \only<12|handout:0> {\includegraphics[width=0.60\linewidth]{\figureLocation{04tours/BoonThaVorn/BTV_Inside09.jpg}}}
    \only<13|handout:0> {\includegraphics[width=0.40\linewidth]{\figureLocation{04tours/BoonThaVorn/BTV_Inside10.jpg}}}
    \only<14|handout:0> {\includegraphics[width=0.40\linewidth]{\figureLocation{04tours/BoonThaVorn/BTV_Inside11.jpg}}}
    \only<15|handout:0> {\includegraphics[width=0.35\linewidth]{\figureLocation{04tours/BoonThaVorn/BTV_Equipment01.jpg}}}
    \only<16|handout:0> {\includegraphics[width=0.35\linewidth]{\figureLocation{04tours/BoonThaVorn/BTV_Equipment02.jpg}}}
    \only<17|handout:0> {\includegraphics[width=0.40\linewidth]{\figureLocation{04tours/BoonThaVorn/BTV_Problem01.jpg}}}
  \end{center}

  \note{
    \noindent\textbf{Purpose:} Connect Warehouse Intuition: Boon Tha Vorn to a decision and the evidence needed to explain it.\par
    \noindent\textbf{Leading Question:} Which design choice supports the Boonthavorn operation?\par
    \noindent\textbf{Teaching cue:} Trace inventory and shipment activities and distinguish evidence from assumptions.\par
    \begin{itemize}
      \item \textbf{Location:} DC near BU (keep SKU type A \& type B) for serving every outlets \& direct shipping customers
      \item \textbf{Business:} retailer of home improvement material (ceramic, sanitary-ware, bathroom fitting)
      \item \textbf{Equipments:} selective rack (4.5T bot, 3.0T mid, 1.5T top), forklift, reach truck, RF barcode scanner
      \item \textbf{Customer service:} .
      \item \textbf{Trivial:} IE Chula 36; investment in WH 200+ bath
    \end{itemize}
  }
  \end{frame}


%================================================================================================-
\subsection[KPI]{Key Performance Index}
%============================================================================================================--
\begin{frame}
  \frametitle{What is KPI?}

  \begin{itemize}
    \item \textbf{What:} a way to measure performance of organization/activity
    \item \textbf{Important:}
    \begin{itemize}
      \item indicate success of each activity
      \item evaluate main objectives
      \item measure progress of implementation (historical comparison)
      \item measure productivity \& efficiency  \NextHalfLine
    \end{itemize}
    \item \textbf{Issues:} data collection, measurement, \alert{consistency}
  \end{itemize}

  \NextLine
  \textbf{Type of KPI?}
  \begin{itemize}
    \item \textbf{Financial related KPI:} \% warehousing cost per total cost, cost per shipped SKU
    \item \textbf{Non-Financial related KPI:}
   \end{itemize}

  \note{
    \noindent\textbf{Purpose:} Connect What is KPI to a decision and the evidence needed to explain it.\par
    \noindent\textbf{Leading Question:} What makes a KPI comparable across warehouses?\par
    \noindent\textbf{Teaching cue:} Specify a definition, unit, sample period, and scope before comparing results.\par
    \begin{itemize}
      \item $<<$empty$>>$
    \end{itemize}
  }
  \end{frame}

%%============================================================================================================--
\begin{frame}
  \frametitle{Example of Non-Financial KPI}

  \begin{itemize}
    \item \textbf{Service Customer View:} \alert{response time} (order cycle time), shipment accuracy (correct qty/total qty), \alert{fill rate} (qty shipped/ordered qty)
    \item \textbf{Service Warehouse View:} dock-to-stock time, \alert{inventory accuracy}, \% cross-docking order
    \item \textbf{Productivity:} \alert{lines per man-hour}, cases per person-hour, \alert{cubic space utilization}, equipment up-time
    \item \textbf{Situation:} lines shipped per SKU, \alert{inventory turnover}, investment pick accuracy, \% of new SKUs, \% active SKUs, labor turnover, lines per order, total lines shipped per day
  \end{itemize}

  \begin{flushright}
    \tiny{\textbf{Adopted from} Hackman, S. 1982. \hskip10em}\\
  \end{flushright}

  \note{
    \noindent\textbf{Purpose:} Connect Example of Non-Financial KPI to a decision and the evidence needed to explain it.\par
    \noindent\textbf{Leading Question:} How could a productivity KPI improve while accuracy deteriorates?\par
    \noindent\textbf{Teaching cue:} Pair an output rate with an order accuracy or damage indicator.\par
    \begin{itemize}
      \item $<<$empty$>>$
    \end{itemize}
  }
  \end{frame}


%============================================================================================================--
\begin{frame}
  \frametitle{Rating of Selected KPIs}


  \begin{center}\small{
  \begin{tabular}{lccccc}\hline
    KPI/Rating 	                 &Poor 	  &Sub-Par 	&Par 	&Superior 	&Outstanding \\ \hline
    responding time (hrs) 	     &$>$48 	&24-48 	  &12-24	  &4-12	   &$<$4  \\
    dock-to-store time (hrs) 	   &$>$48 	&24-48 	  &8-24 	  &2-8	   &$<$2 \\
    lines per man-hour 	         &$<$5 	  &5-10	    &10-20	  &20-50 	 &$>$50 \\
    cases per man-hour 	         &8-25	  &25-50 	  &50-100 	&100-250 	&$>$250 \\
    cubic utilization (\%) 	     &$<$65   &65-75 	  &75-85 	  &85-95 & $>$95 \\
    annual lines per SKU 	       &$<$50 	&50-100	  &100-250 	&250-400	&$>$400 \\
    inventory discrepancy (\% qty)  &$>$5.0 	&1.0-5.0 	&0.5-1.0 	&0.05-0.5	 &$<$0.05 \\
    inventory turn 	             &$<$1.0 	&1.0-3.0 	&3.0-6.0 	&6.0-10.0  &$>$12.0 \\ \hline
  \end{tabular} }
  \end{center}

  \begin{flushright}
    \tiny{\textbf{source} Hackman, S. \emph{et al.} 2001~\cite{hackman2001benchmarking,}. \hskip10em}\\
  \end{flushright}


  \note{
    \noindent\textbf{Purpose:} Connect Rating of Selected KPIs to a decision and the evidence needed to explain it.\par
    \noindent\textbf{Leading Question:} Which of these benchmarks need unit or direction checks before use?\par
    \noindent\textbf{Teaching cue:} Read 0.95 as 95 percent in a percentage row and verify the original bands and discrepancy definition.\par
    \begin{itemize}
      \item $<<$empty$>>$
    \end{itemize}
  }

  \end{frame}

%============================================================================================================--
\section[Challenges]{Typical Challenges in Warehouse}
%============================================================================================================--

  \begin{frame}
  \frametitle{Nature of warehouse}

  \begin{itemize}
    \item Warehouse is \alert{labor intensive}
    \item Warehouse is, in general, the \alert{last frontier} in Supply Chain before start collaborate
    \item Investment in warehouse depends on \alert{values of SKUs} in warehouse
    \item \alert{Broken pallets/cases} tend to be damaged \& lost
    \item Flows of material in a warehouse is \alert{rarely balance} at particular time
    \item Works \& effort warehousing activities are \alert{unbalance}
    \begin{itemize}
      \item<2-> Put-Away $\leq$ Pick-Up
      \item<2-> Receiving $\leq$ Shipping
    \end{itemize}
  \end{itemize}

  \note{
    \noindent\textbf{Purpose:} Connect Nature of warehouse to a decision and the evidence needed to explain it.\par
    \noindent\textbf{Leading Question:} Which structural features make one warehouse different from another?\par
    \noindent\textbf{Teaching cue:} Compare product mix, order profile, scale, and delivery promise.\par
    \begin{itemize}
      \item \textbf{labor intensive:} need labor to pick/pack unless volume very high
      \item \textbf{values:} gypsum board VS ceramic tile or Raw mat' VS FG
      \item \textbf{broken pallet:} dilemma of picking 8/10 pens
      \item \textbf{balance:} this about X-mas
    \end{itemize}
  }
  \end{frame}


%============================================================================================================
\usetikzlibrary{arrows,positioning,decorations.pathmorphing,calc}
\renewcommand{\textFont}[1]{\footnotesize{#1}}
%==============================================================================================================
\begin{frame}
  \frametitle{Typical issues in warehouse}

  \begin{center}
  \begin{tikzpicture}[xscale=1.20,yscale=1.20,>=stealth]
    \tikzstyle{block}=[rectangle,draw=black!80,thick,fill=blue!15,text centered,text width=30mm,text height=4mm];

    \node[]       (ref) {\textFont{~}};
    \node[block]  (nospace)   [left = of ref,xshift=-15mm]  {\textbf{No space}} ;
    \node[block]  (inaccuate) [right = of ref,xshift=15mm]  {\textbf{Inaccurate}} ;
    \node[block]  (prod)      [below = of ref,yshift=-15mm] {\textbf{Productivity}} ;

    \draw[<->,thick] (nospace) --node[align=left,below] {\textFont{
    \begin{tabular}{l}
        $\bullet$ personal feeling \\
        $\bullet$ disciple of worker
      \end{tabular}
    }}  (inaccuate) ;
    \draw[<->,thick] (nospace) --node[align=left,below] {\textFont{
    \begin{tabular}{l}
        $\bullet$ manual update/ searching \\
        $\bullet$ temp bin location
      \end{tabular}
    }}  (prod) ;
    \draw[<->,thick] (prod) --node[align=left,below right] {\textFont{
    \begin{tabular}{l}
        $\bullet$ lack worker \\
        $\bullet$  complied to std \\
        $\bullet$ sheer transactions
      \end{tabular}
    }} (inaccuate) ;

    \node[above,align=left] at (nospace.north) { \textFont{ \code{
    \begin{tabular}{l}
        $\bullet$ slotting \\
        $\bullet$ \# SKU \\
        $\bullet$ excess qty
      \end{tabular} }} };

    \node[above,align=left] at (inaccuate.north) { \textFont{ \code{
    \begin{tabular}{l}
        $\bullet$ freq counting \\
        $\bullet$ careless receiving/ picking
      \end{tabular} }} };

    \node[below,align=left] at (prod.south) { \textFont{ \code{
    \begin{tabular}{l}
        $\bullet$ lack workers \\
        $\bullet$ economy of scale $\rightarrow$ equipment
      \end{tabular} }} };

  \end{tikzpicture}
  \end{center}

  \begin{itemize}
      \item \textbf{Other:} fire, safety, quality, use of data
  \end{itemize}

  \note{
    \noindent\textbf{Purpose:} Connect Typical issues in warehouse to a decision and the evidence needed to explain it.\par
    \noindent\textbf{Leading Question:} What evidence would locate the cause of a recurring warehouse problem?\par
    \noindent\textbf{Teaching cue:} Observe one process, record the exception rate, and separate symptoms from causes.\par
    \begin{itemize}
      \item $<<$empty$>>$
    \end{itemize}
  }
  \end{frame}

%================================================================================================-
\begin{frame}
  \frametitle{Summary: Warehouse $\in$ Supply Chain}

  \only<1>{
  \begin{itemize}
    \item \textbf{Warehouse} $\neq$ \alert{storage room} because balance flows, system, policy, equipments, \emph{etc}
    \item \textbf{Warehousing management:} maximizing usage of warehouse 'resources' at 'right' service level
    \item \textbf{Warehousing activities:} focus on \alert{main} activities, minimize \alert{non-value} activities
    \item \textbf{Equipments:} easy to handling/ storage/ track;
    \item \textbf{Layouts:} put a \alert{fast moving} items in \alert{convenient locations}
  \end{itemize}
  }

  \only<1>{
    \NextHalfLine
    \begin{alertblock}{}{}{}
      \textbf{F.A.S.T} = \alert{F}low/ \alert{A}ccessibility/ \alert{S}pace/ \alert{T}hroughput
    \end{alertblock}
  }

  \only<2|handout:0>{
  \begin{center}
    \includegraphics[width=0.70\textwidth]{\figureLocation{04layout/FAST-Flow00.jpg}} \\
    \textbf{Ideal:} Uninterrupted \& 'shortest' \alert{Flow}; like \alert{\textbf{water system}}
  \end{center}
  }

  \only<3|handout:0>{
  \begin{center}
  \begin{tabular}{cc}
    \begin{minipage}{0.45\linewidth}
      \includegraphics[width=0.65\textwidth]{\figureLocation{04layout/FAST-Access01.jpg}}
    \end{minipage}
    &
    \begin{minipage}{0.45\linewidth}
      \includegraphics[width=0.95\textwidth]{\figureLocation{04layout/FAST-Access02.jpg}}
    \end{minipage}
  \end{tabular}
    \NextHalfLine
  \textbf{Ideal:} Easy \alert{Accessibility} in \& out; congestion $=$ \alert{\textbf{money}}
  \end{center}
  }

  \only<4|handout:0>{
  \begin{center}
  \begin{tabular}{cc}
    \begin{minipage}{0.45\linewidth}
      \includegraphics[width=0.95\textwidth]{\figureLocation{04layout/FAST-Space01.jpg}}
    \end{minipage}
    &
    \begin{minipage}{0.45\linewidth}
      \includegraphics[width=0.95\textwidth]{\figureLocation{04layout/FAST-Space02.jpg}}
    \end{minipage}
  \end{tabular}
    \NextHalfLine
  \textbf{Ideal:} Maximizing \alert{Space} unitization; space is \alert{\textbf{not free}}
  \end{center}
  }

  \only<5|handout:0>{
    \begin{center} %\small{
    \begin{tabular}{lr} \hline
      Throughput 	              & Value \\ \hline
      No. of SKUs on inventory 	& 60,000 \\
      SKUs stored in DC 	      & 10,000 \\
      SKUs cross-docked 	      & 50,000 \\
      Average stock holding  	  & 12 days\\
      Average cases/pallet      & 50 \\
      Average pieces/case 	    & 24 \\
      Non-Std pieces/pallet     & 4 \\
      Returned goods   	        & 5\% \\ \hline
    \end{tabular} %}
    \NextHalfLine
    \textbf{Ideal:} Maximizing \alert{thoughtful}, warehouse \alert{\textbf{$\neq$ storage room}}
    \end{center}
  }


  \note{
    \noindent\textbf{Purpose:} Connect Summary: Warehouse ∈ Supply Chain to a decision and the evidence needed to explain it.\par
    \noindent\textbf{Leading Question:} Which warehouse choice affects performance elsewhere in the supply chain?\par
    \noindent\textbf{Teaching cue:} Explain the impact of one decision on service, inventory, or transport outside the site.\par
    \begin{itemize}
    \item Warehousing is related to:
    \begin{itemize}
        \item<2-> \textbf{Direct:}   inventory control, material handling, distribution network, packaging
        \item<2-> \textbf{Indirect:} facility design, assembly, production, purchasing, customer service
    \end{itemize}
    \begin{itemize}
      \item why storage is not activity in WH (most time consuming in WH)?
      \item 6Right in Supply Chain (Item, Qty, Condition, Time, Place, Price)
    \end{itemize}
    \end{itemize}
  }
\end{frame}

%============================================================================================================--
\begin{frame}
  \frametitle{WH WHS MGT Textbook}

 \begin{center}
  \begin{minipage}{0.75\textwidth}
    \begin{tabular}{ccc}
      \includegraphics[width=0.45\textwidth]{\figureLocation{coverBook.png}} &
      \includegraphics[width=0.45\textwidth]{\figureLocation{snapShot.png}} \\
    \end{tabular}
  \end{minipage}
 \end{center}

  \small{
  \begin{itemize}
    \item \textbf{Publisher:}  G.P Cyber Print~\cite{oran2018warehouseENG}
    \item \textbf{ISBN:} 978-6164073722
  \end{itemize}
  }

  \note{
    \noindent\textbf{Purpose:} Connect WH WHS MGT Textbook to a decision and the evidence needed to explain it.\par
    \noindent\textbf{Leading Question:} Which reference chapter could help validate your layout or KPI choice?\par
    \noindent\textbf{Teaching cue:} Identify one topic and cross-check the definition against the displayed source.\par
    a

  }
  \end{frame}

%================================================================================================-

\begin{frame}
  \frametitle{Reference}

  \begin{footnotesize}
    \bibliography{\refLocation{Warehouse}}
  \end{footnotesize}

  \note{
    \noindent\textbf{Purpose:} Help students verify and reuse evidence from the lecture.\par
    \noindent\textbf{Leading Question:} What would you check before reusing a figure, fact, or model from this lecture?\par
    \noindent\textbf{Teaching cue:} Trace a citation to its original source and verify the year, scope, and assumptions.\par
  }

\end{frame}


%================================================================================================-
\end{document}

\end{frame} 